\chapter{Metodologia}\label{chap:metodologia}
This chapter describes the tools used and the tests performed for this work. The uplink transmission of a massive MIMO system with decentralized processing (Section~\ref{subsec:uplink}) was simulated to analyze the effect of the approximation used in the inversion module of the BCD-MMSE algorithm. Based on this analysis, the fixed-point formats for the inversion module, implemented on an FPGA, are selected. The chapter begins by presenting the simulation flow (Section~\ref{sec:simulador}) and the hardware tools (Section~\ref{sec:setup-hardware}). It then covers test configurations and conditions (Section~\ref{sec:configuracoes}), timing constraints for the inversion process (Section~\ref{sec:requisitos-tempo}), the selection of the numerical format (Section~\ref{sec:formato}), the hardware architectures and testing (Section~\ref{sec:testes-hardware}), and the organization of the results (Section~\ref{sec:organizacao-resultados}).



% =====================================================================
\section{\texorpdfstring{Simulation workflow}{Simulation workflow}}\label{sec:simulador}
% =====================================================================
The simulation workflow~\ref{fig:fluxo_simulacao} employs three tools: the QuaDRiGa channel generator, a Python simulator developed in this work, and C++ code developed in Vitis HLS, which computes the inverse of $\mathbf{G}_c$ using fixed-point arithmetic. QuaDRiGa generates the channel matrices $\mathbf{H}$ for the target UEs and $\overline{\mathbf{H}}$ for the interfering UEs.

The Python simulator handles the transmission process. It generates symbols for the target UEs and creates colored noise, with a covariance derived from $\overline{\mathbf{H}}$ (Section~\ref{subsec:ruido-colorido}), to construct the received signal using $\mathbf{H}$. Subsequently, it computes the equalizers, performs symbol detection, and measures the SER at various SNR levels. The BCD-MMSE algorithm is simulated in two ways: using the inverse of $\mathbf{G}_c$ in 64-bit floating-point format (which serves as a reference) and using the inverse in fixed-point format. In the latter case, whenever the BCD-MMSE algorithm requires $\mathbf{G}_c^{-1}$, the Python simulator invokes the C++ code.

The Python simulator models the same uplink as in \cite{zhao2025jsac}. According to the authors, that paper was the first work on decentralized processing to consider colored noise \cite{zhao2025jsac}. Therefore, the simulation parameters follow that work, except where the text indicates otherwise.



% Fluxo conferido no código:
%   QuaDRiGa: software/simulador/quadriga/export_quadriga_channels.m (semente: linha 83)
%   laço e equalizadores: dbp_sim/m1_driver.py:149-211; dbp_sim/m7_equalizers.py
%     (lmmse, zero_forcing, bdac_mmse, bcd_mmse_path); Ledoit-Wolf: dbp_sim/m5_noise.py:162-168
%   inversa em ponto fixo: dbp_sim/hls_inverse.py (usada nos resultados de SER,
%     varredura_formatos.py, bcd_inverse="hls") e dbp_sim/hw_inverse.py (modelo em Python)
% Conteúdo vem de constraints/diagramas/fluxo_simulador.puml; estrutura em caixas do
% rascunho constraints/diagramas/metodologia.svg.
\begin{figure}[!ht]
\caption{Simulation flow: QuaDRiGa channels, Python simulator and fixed-point inverse.}
\label{fig:fluxo_simulacao}
\vspace{4pt} % a classe zera o espaço entre a caption e a figura
\centering
\begin{tikzpicture}[
  font=\fontsize{7}{8.4}\selectfont,
  >={Stealth[length=2.1mm, width=1.7mm]},
  node distance=2.4mm,
  box/.style={draw=black!65, fill=white, align=center,
              inner sep=2.5pt, text width=#1,
              execute at begin node={\hyphenpenalty=10000}},
  box/.default=4.9cm,
  inp/.style={trapezium, trapezium left angle=78, trapezium right angle=102,
              trapezium stretches body, draw=black!65, fill=black!6,
              align=left, inner xsep=9pt, inner ysep=3pt},
  arq/.style={tape, tape bend top=none, draw=black!65, fill=white, align=center,
              text width=#1, inner sep=2.5pt},
  pre/.style={box=#1, inner xsep=2.5pt+1.3mm,
              append after command={
                ($(\tikzlastnode.north west)+(1.1mm,0)$) edge[draw=black!65, -]
                ($(\tikzlastnode.south west)+(1.1mm,0)$)
                ($(\tikzlastnode.north east)+(-1.1mm,0)$) edge[draw=black!65, -]
                ($(\tikzlastnode.south east)+(-1.1mm,0)$)}},
  lini/.style={chamfered rectangle, chamfered rectangle corners={north west, north east},
              chamfered rectangle xsep=2mm, chamfered rectangle ysep=0.6mm,
              draw=black!65, fill=black!6, align=center, inner sep=2.5pt, text width=#1},
  lini/.default=4.2cm,
  lfim/.style={lini=#1, chamfered rectangle corners={south west, south east}},
  lfim/.default=4.2cm,
  term/.style={rounded rectangle, draw=black!65, fill=black!6, align=center,
              inner xsep=8pt, inner ysep=3pt},
  dec/.style={diamond, aspect=3.2, draw=black!65, fill=black!6, align=center, inner sep=0.5pt},
  ponto/.style={circle, draw=black!65, fill=white, minimum size=2.2mm, inner sep=0pt},
  cont/.style={draw=black!65, dashed, fill=black!4},
  tit/.style={anchor=north west, font=\fontsize{8}{9.6}\selectfont\bfseries\boldmath, align=left, inner sep=3pt},
  seta/.style={->, draw=black, line width=0.6pt},
  rot/.style={font=\fontsize{5.5}{6.6}\selectfont, text=black!70, inner sep=1.5pt},
]

% Colunas: QuaDRiGa [0; 5,4], Python [5,7; 9,0] (mais estreita, só um bloco),
%   inversa [9,3; 14,7]; 0,3 cm entre elas. Caixas tracejadas com a mesma altura;
%   entradas com o mesmo tamanho e alinhadas; setas das entradas à direita do centro
%   de cada coluna (fora dos títulos).
% =====================================================================
% Entradas
% =====================================================================
\begin{scope}[ent/.style={inp, anchor=north, text width=3.0cm, minimum height=1.5cm,
                           execute at begin node={\hyphenpenalty=10000\raggedright}}]
\node[term] (ini) at (2.7,0.75) {Start};
\node[ent] (inq) at (2.7,0) {\textbf{Inputs:} $M$, $K$, $r$, carrier frequency, propagation scenario, seed};
\node[ent] (ins) at (7.35,0) {\textbf{Inputs:} configuration, test condition, SNR points, realizations, symbols per UE, BCD-MMSE iterations};
\node[ent] (inf) at (12.0,0) {\textbf{Inputs:} format $\langle W, I\rangle$, quantization and overflow modes};
\end{scope}

% topo do conteúdo das colunas (as caixas tracejadas começam 9 mm acima)
\path let \p1=(ins.south) in coordinate (topo) at (7.35,\y1-14mm);

% =====================================================================
% QuaDRiGa (coluna da esquerda)
% =====================================================================
\path let \p1=(topo) in
  node[box=3.95cm, anchor=north] (exp) at (2.7,\y1) {\texttt{export\_quadriga\_channels.m}\\ once per configuration $(M, K, r)$\\ and carrier};
\node[lini=3.3cm, below=of exp] (lq) {for each realization $i$ (100)};
\node[box=3.3cm, below=of lq] (rea) {\mbox{seed $= i$}; draws the UE positions, generates $\mathbf{H}$ for the target UEs and $\overline{\mathbf{H}}$ for the interfering UEs};
\node[lfim=3.3cm, below=of rea] (lqf) {next realization};
\node[arq=2.8cm, below=of lqf] (mat) {all channel matrices\\ in one \texttt{.mat} file};
\node[pre=3.95cm, below=of mat] (cnv) {\texttt{convert\_quadriga\_export.py}\\ \texttt{.mat} to \texttt{.npz} (NumPy file)};
\node[arq=2.8cm, below=of cnv] (npz) {all channel matrices\\ in one \texttt{.npz} file};

\draw[seta] (ini) -- (inq);
\draw[seta] let \p1=(inq.south), \p2=(exp.north) in (4.3,\y1) -- (4.3,\y2);
\draw[seta] (exp) -- (lq);
\draw[seta] (lq) -- (rea);
\draw[seta] (rea) -- (lqf);
\draw[seta] (lqf) -- (mat);
\draw[seta] (mat) -- (cnv);
\draw[seta] (cnv) -- (npz);

% =====================================================================
% Inversa em ponto fixo (coluna da direita)
% =====================================================================
\path let \p1=(topo) in
  node[dec, aspect=2.4] (bcdi) at (12.0,\y1-1.0cm) {\texttt{bcd\_inverse}};
\path let \p1=(bcdi.south) in
  node[pre=1.9cm, anchor=north] (hls) at (13.4,\y1-5mm) {{\fontsize{6.2}{7.4}\selectfont\texttt{hls\_inverse.py}}\\ sends $\mathbf{G}_c$\\ through a pipe};
\path let \p1=(bcdi.south) in
  node[pre=1.7cm, anchor=north] (gold) at (10.6,\y1-5mm) {{\fontsize{6.2}{7.4}\selectfont\texttt{hw\_inverse.py}}\\ bit-exact Python model (golden model), checked against the C-simulation};
\node[box=1.95cm, below=7mm of hls] (cfx) {computes the inverse in fixed point with the compiled HLS C++ code};
\node[box=1.95cm, below=of cfx] (dev) {returns the inverse through the pipe};
\node[rot, align=center, below=9.5mm of dev] (uso) {\texttt{"hls"} was used\\ for the SER results};

\draw[seta] let \p1=(inf.south), \p2=(bcdi.north) in (13.6,\y1) -- (13.6,\y2+3mm) -- (\x2,\y2+3mm) -- (\p2);
\draw[seta] (bcdi.east) -| node[rot, pos=0.8, left] {\texttt{"hls"}} (hls.north);
\draw[seta] (bcdi.west) -| node[rot, pos=0.8, right] {\texttt{"cholesky\_hw"}} (gold.north);
\draw[seta] (hls) -- (cfx);
\draw[seta] (cfx) -- (dev);

% fundo comum das três caixas tracejadas
\path let \p1=(npz.south), \p2=(uso.south), \n1={min(\y1,\y2)-2.5mm} in coordinate (fundo) at (7.35,\n1);

% =====================================================================
% Simulador em Python: um bloco (detalhado na Figura fig:fluxo_python)
% =====================================================================
\path let \p1=(fundo) in
  node[arq=2.3cm, anchor=south] (crv) at (7.2,\y1+2.5mm) {SER $\times$ SNR curves};
\path let \p1=(topo), \p2=(crv.north), \n1={\y1-\y2-3mm} in
  node[box=2.1cm, anchor=north, minimum height=\n1] (py) at (7.2,\y1) {Monte Carlo\\ simulation\\ of the uplink\\[2pt] (detailed in\\ Figure~\ref{fig:fluxo_python})};
\node[term, below=7mm of crv] (fim) {End};

\draw[seta] let \p1=(ins.south), \p2=(py.north) in (8.1,\y1) -- (8.1,\y2);
\draw[seta] (py) -- (crv);
\draw[seta] (crv) -- (fim);

% canais do QuaDRiGa para o simulador
\draw[seta] (npz.east) -- (npz.east -| py.west);

% inversa: simulador (BCD-MMSE update) <-> coluna da direita
\draw[seta] let \p1=(py.east), \p2=(hls) in (\x1,\y2+1.5mm) -- node[rot, above, fill=white, inner sep=0.6pt, pos=0.5] {$\mathbf{G}_c$} (9.3,\y2+1.5mm);
\draw[seta] let \p1=(py.east), \p2=(hls) in (9.3,\y2-1.5mm) -- node[rot, below, fill=white, inner sep=0.6pt, pos=0.5] {$\mathbf{G}_c^{-1}$} (\x1,\y2-1.5mm);

% =====================================================================
% Caixas tracejadas (fundo)
% =====================================================================
\begin{scope}[on background layer]
  \path let \p1=(topo) in coordinate (qt) at (0,\y1+9mm);
  \draw[cont] (qt) rectangle (5.4,0 |- fundo);
  \node[tit] at (qt) {QuaDRiGa 2.8.1 (MATLAB)};

  \path let \p1=(topo) in coordinate (rt) at (5.7,\y1+9mm);
  \draw[cont] (rt) rectangle (9.0,0 |- fundo);
  \node[tit] at (rt) {Python simulator\\ (\texttt{dbp\_sim})};

  \path let \p1=(topo) in coordinate (it) at (9.3,\y1+9mm);
  \draw[cont] (it) rectangle (14.7,0 |- fundo);
  \node[tit] at (it) {Fixed-point inverse\\ of $\mathbf{G}_c$ ($M_c \times M_c$)};

  \path let \p1=(cfx.north), \p2=(dev.south) in
    coordinate (vt) at (12.2,\y1+2mm) coordinate (vb) at (14.55,\y2-8.5mm);
  \draw[cont, fill=black!9] (vt) rectangle (vb);
  \node[tit, font=\fontsize{7}{8.4}\selectfont\bfseries, anchor=south west] at (vt |- vb) {g++ with Vitis\\ HLS headers};
\end{scope}

\end{tikzpicture}
\vspace{4pt}
\legend{Source: the author.}
\end{figure}

% ---------------------------------------------------------------------
\subsection{\texorpdfstring{Channel generation}{Channel generation}}\label{subsec:simul-canal}
% ---------------------------------------------------------------------
The channels are generated using QuaDRiGa 2.8.1 in MATLAB. They depend on the number of antennas $M$, the number of target UEs $K$, the number of interfering UEs $r$, and the carrier frequency. Consequently, they are generated once for each combination of these values. Other parameters, such as $C$, the modulation scheme, and the IoTh, only affect the processing steps within the Python-based simulator. The simulation employs 100 channel realizations. Each realization consists of a random selection of UE positions, along with the corresponding channels $\mathbf{H}$ and $\overline{\mathbf{H}}$ generated by QuaDRiGa for those specific positions. Every realization uses its own unique random seed.


MATLAB saves the channels for all 100 realizations into a \texttt{.mat} file. A Python script then converts this file into an \texttt{.npz} file, NumPy’s format for storing multiple arrays within a single compressed file. The Python-based simulator reads the matrices $\mathbf{H}$ and $\overline{\mathbf{H}}$ from this file.
% ---------------------------------------------------------------------
\subsection{\texorpdfstring{Python simulator}{Python simulator}}\label{subsubsec:simul-python}
% ---------------------------------------------------------------------
The Python simulator runs two nested loops (Figure~\ref{fig:fluxo_python}). The outer loop iterates through the 100 realizations, while the inner loop iterates through the SNR points. The SNR points range from $-10$ to $10$~dB, with a step of $2$~dB. For each realization, the simulator takes $\mathbf{H}$ and $\overline{\mathbf{H}}$ corresponding to that realization and uses these channels across all SNR points. Each point on the SER curves represents the average over the 100 realizations. This is a Monte Carlo simulation, as in \cite{zhuang2025}: the SER (Section~\ref{subsubsec:ser}) is estimated by the fraction of erroneous symbols across many random draws of the channel, noise, and symbols.

At each SNR point, the simulator constructs the noise covariance matrix based on $\overline{\mathbf{H}}$, the IoTh (interference-over-thermal), and the SNR. It then generates $N$ noise samples, $\mathbf{n}^i$ (Section~\ref{subsubsec:estimacao-R}). It also generates 480 data symbols for each UE, along with the noise associated with these symbols. Based on this, the equalizers are calculated. The channel remains constant within the coherence block (Section~\ref{subsec:bloco-coerencia}); therefore, each equalizer is computed once and applied to all 480 symbols. The BCD-MMSE is compared against two centralized equalizers, the LMMSE and the ZF, which serve as benchmarks. The ZF does not utilize the noise samples $\mathbf{n}^i$. In contrast, both the LMMSE and the BCD-MMSE employ these samples to estimate the noise covariance (Section~\ref{subsubsec:condicionamento-R}). This estimate is always subjected to Ledoit–Wolf shrinkage (Equation~\ref{eq:ledoit-wolf}) before being used, even in cases where the estimate is already invertible.

In the LMMSE, this adjustment is applied to the estimate $\hat{\mathbf{R}}$ corresponding to all $M$ antennas. In the BCD-MMSE, each DU performs the adjustment on its own block, $\hat{\mathbf{R}}_{cc}$, using only the samples from its $M_c$ antennas. Consequently, no DU requires access to the samples from any other DU. The adjusted block forms the matrix $\mathbf{G}_c$ (Section~\ref{subsec:bcd}). In the reference version of the BCD-MMSE, this matrix is inverted using 64-bit floating-point arithmetic; alternatively, it can be inverted using fixed-point arithmetic (Section~\ref{subsubsec:simul-inversa}).



\begin{figure}[!ht]
\caption{Detailed flow of the Python simulator: Monte Carlo loop, equalizers and SER.}
\label{fig:fluxo_python}
\vspace{4pt} % a classe zera o espaço entre a caption e a figura
\centering
\begin{tikzpicture}[
  font=\fontsize{7}{8.4}\selectfont,
  >={Stealth[length=2.1mm, width=1.7mm]},
  node distance=2.4mm,
  box/.style={draw=black!65, fill=white, align=center,
              inner sep=2.5pt, text width=#1,
              execute at begin node={\hyphenpenalty=10000}},
  box/.default=4.9cm,
  inp/.style={trapezium, trapezium left angle=78, trapezium right angle=102,
              trapezium stretches body, draw=black!65, fill=black!6,
              align=left, inner xsep=9pt, inner ysep=3pt},
  arq/.style={tape, tape bend top=none, draw=black!65, fill=white, align=center,
              text width=#1, inner sep=2.5pt},
  pre/.style={box=#1, inner xsep=2.5pt+1.3mm,
              append after command={
                ($(\tikzlastnode.north west)+(1.1mm,0)$) edge[draw=black!65, -]
                ($(\tikzlastnode.south west)+(1.1mm,0)$)
                ($(\tikzlastnode.north east)+(-1.1mm,0)$) edge[draw=black!65, -]
                ($(\tikzlastnode.south east)+(-1.1mm,0)$)}},
  lini/.style={chamfered rectangle, chamfered rectangle corners={north west, north east},
              chamfered rectangle xsep=2mm, chamfered rectangle ysep=0.6mm,
              draw=black!65, fill=black!6, align=center, inner sep=2.5pt, text width=#1},
  lini/.default=4.2cm,
  lfim/.style={lini=#1, chamfered rectangle corners={south west, south east}},
  lfim/.default=4.2cm,
  term/.style={rounded rectangle, draw=black!65, fill=black!6, align=center,
              inner xsep=8pt, inner ysep=3pt},
  dec/.style={diamond, aspect=3.2, draw=black!65, fill=black!6, align=center, inner sep=0.5pt},
  ponto/.style={circle, draw=black!65, fill=white, minimum size=2.2mm, inner sep=0pt},
  cont/.style={draw=black!65, dashed, fill=black!4},
  tit/.style={anchor=north west, font=\fontsize{8}{9.6}\selectfont\bfseries\boldmath, align=left, inner sep=3pt},
  seta/.style={->, draw=black, line width=0.6pt},
  rot/.style={font=\fontsize{5.5}{6.6}\selectfont, text=black!70, inner sep=1.5pt},
  lp/.style={draw=black!65, rounded corners=3pt, fill=#1},
  lab/.style={anchor=north east, font=\fontsize{7}{8.4}\selectfont\bfseries, inner sep=3pt},
  grp/.style={font=\fontsize{6.5}{7.8}\selectfont\bfseries, text=black!70, inner sep=1pt},
  bus/.style={draw=black, line width=0.6pt},
]

% Organização (2026-11-05): entrada em coluna acima de Signal and noise, laços como molduras
%   (realizações por fora, pontos de SNR por dentro), três grupos lado a lado dentro do
%   laço de SNR (Signal and noise | Receiver | Evaluation, os grupos das tabelas de
%   Resultados) e saída em coluna abaixo de Evaluation. Conteúdo das caixas igual ao da
%   versão anterior (cópia: scratchpad/simulacao_antes_detalhada.tex).

% =====================================================================
% Entrada: coluna vertical acima do grupo Signal and noise (x = 1,85)
% =====================================================================
\node[term] (ini) at (1.85,0) {Start};
\node[inp, anchor=north, text width=4.4cm, execute at begin node={\hyphenpenalty=10000\raggedright}] (ins) at (3.0,-0.55) {%
  \textbf{Inputs:} configuration ($M$, $C$, $K$, IoTh, modulation), test condition ($N$, carrier frequency), SNR points, realizations, symbols per UE, BCD-MMSE iterations};
\path let \p1=(ins.south) in
  node[box=3.2cm, anchor=north] (rd) at (1.85,\y1-3mm) {reads the channels $\mathbf{H}$ and $\overline{\mathbf{H}}$ (\texttt{.npz})};
\node[box=3.2cm, below=3mm of rd] (gen) {starts a random generator with \mbox{seed $= 0$}, once per run; it draws the noise samples, the symbols and the noise};
\draw[seta] (ini) -- (ini |- ins.north);
\draw[seta] (rd.north |- ins.south) -- (rd);
\draw[seta] (rd) -- (gen);

% referências verticais
\path let \p1=(gen.south) in
  coordinate (yo) at (0,\y1-4mm)       % topo da moldura das realizações
  coordinate (yi) at (0,\y1-10mm)      % topo da moldura dos pontos de SNR
  coordinate (y0) at (0,\y1-16mm);     % topo do conteúdo

% =====================================================================
% Receiver: três faixas (BCD-MMSE, ZF, LMMSE)
% =====================================================================
\path let \p1=(y0) in
  node[box=2.1cm, anchor=north] (d4) at (10.35,\y1) {BCD-MMSE update in each DU: $\mathbf{G}_c = E_s \mathbf{H}_c \mathbf{H}_c^H$ $+$ covariance of DU $c$; $\mathbf{G}_c^{-1}$ in float64 or fixed point (Figure~\ref{fig:fluxo_simulacao})};
\path let \p1=(d4) in
  node[box=2.1cm] (d1) at (4.95,\y1) {\textbf{Decentralized estimation:} each DU $c$ makes a Ledoit-Wolf estimate of the covariance of the $M_c$ entries of $\mathbf{n}^i$ from its antennas ($M_c \times M_c$)}
  node[box=2.1cm] (d3) at (7.65,\y1) {BDAC-MMSE initialization (float64)};
\path let \p1=(d1.south), \p2=(d4.south), \n1={min(\y1,\y2)} in
  node[box=2.1cm, anchor=north] (z1) at (4.95,\n1-3.5mm) {ZF equalizer (float64); does not use $\mathbf{n}^i$};
\node[box=2.1cm, below=3.5mm of z1] (c1) {\textbf{Centralized estimation:} Ledoit-Wolf estimate of the covariance $\hat{\mathbf{R}}_n$ ($M \times M$)};
\path let \p1=(c1) in node[box=2.1cm] (c3) at (7.65,\y1) {LMMSE equalizer (float64)};
\draw[seta] (d1) -- (d3);
\draw[seta] (d3) -- (d4);
\draw[seta] (c1) -- (c3);

% centro vertical dos grupos laterais
\path let \p1=(d4), \p2=(c1), \n1={(\y1+\y2)/2} in coordinate (ym) at (0,\n1);

% =====================================================================
% Signal and noise (esquerda) e Evaluation (direita)
% =====================================================================
\path let \p1=(ym) in
  node[box=2.3cm, anchor=south] (cov) at (1.85,\y1+1.75mm) {builds the noise covariance from $\overline{\mathbf{H}}$, IoTh and the SNR point}
  node[box=2.3cm, anchor=north] (smp) at (1.85,\y1-1.75mm) {the generator draws the $N$ noise samples $\mathbf{n}^i$ ($M \times 1$ each), used by the LMMSE and the BCD-MMSE};
\draw[seta] (cov) -- (smp);
\path let \p1=(ym) in
  node[box=2.4cm, anchor=south] (sym) at (13.3,\y1+1.75mm) {the generator draws 480 symbols per UE and their noise: $\mathbf{y} = \mathbf{H}\mathbf{s} + \mathbf{n}$}
  node[box=2.4cm, anchor=north] (ser) at (13.3,\y1-1.75mm) {each equalizer detects the symbols; SER of this realization and SNR point $=$ wrong symbols$\,/\,480K$};
\draw[seta] (sym) -- (ser);

% barramentos: amostras -> três equalizadores -> símbolos
\draw[bus] (smp.east) -- (smp.east -| 3.45,0);
\draw[bus] (d1 -| 3.45,0) -- (c1 -| 3.45,0);
\draw[seta] (d1 -| 3.45,0) -- (d1.west);
\draw[seta] (z1 -| 3.45,0) -- (z1.west);
\draw[seta] (c1 -| 3.45,0) -- (c1.west);
\draw[bus] (d4.east) -- (d4 -| 11.75,0);
\draw[bus] (z1.east) -- (z1 -| 11.75,0);
\draw[bus] (c3.east) -- (c3 -| 11.75,0);
\draw[bus] (d4 -| 11.75,0) -- (c3 -| 11.75,0);
\draw[seta] (sym.west -| 11.75,0) -- (sym.west);

% =====================================================================
% Molduras dos laços e nomes dos grupos
% =====================================================================
\path let \p1=(c1.south), \p2=(c3.south), \p3=(ser.south), \p4=(smp.south),
          \n1={min(\y1,\y2,\y3,\y4)} in
  coordinate (yg) at (0,\n1-3mm)       % nomes dos grupos
  coordinate (yib) at (0,\n1-6.5mm)    % fundo da moldura dos pontos de SNR
  coordinate (yob) at (0,\n1-9.5mm);   % fundo da moldura das realizações
\node[grp] at (1.85,0 |- yg) {Signal and noise};
\node[grp] at (7.65,0 |- yg) {Receiver};
\node[grp, anchor=east] at (13.15,0 |- yg) {Evaluation};

\begin{scope}[on background layer]
  \draw[lp=black!3] (0,0 |- yo) rectangle (15.0,0 |- yob);
  \draw[lp=black!8] (0.3,0 |- yi) rectangle (14.7,0 |- yib);
\end{scope}
\node[lab] at (15.0,0 |- yo) {Monte Carlo loop: for each realization (100), with its $\mathbf{H}$ and $\overline{\mathbf{H}}$};
\node[lab] at (14.7,0 |- yi) {for each SNR point (11)};

% entrada nos laços: gerador -> covariância (reta)
\draw[seta] (gen) -- (cov);

% =====================================================================
% Saída: coluna vertical abaixo do grupo Evaluation (x = 13,3)
% =====================================================================
\path let \p1=(yob) in
  node[box=2.6cm, anchor=north] (avg) at (13.3,\y1-4mm) {average of the 100 SERs of each SNR point};
\node[arq=2.6cm, below=3mm of avg] (crv) {SER $\times$ SNR curves};
\node[term, below=4mm of crv] (fim) {End};
\draw[seta] (ser) -- (avg);
\draw[seta] (avg) -- (crv);
\draw[seta] (crv) -- (fim);

\end{tikzpicture}
\vspace{4pt}
\legend{Source: the author.}
\end{figure}

% ---------------------------------------------------------------------
\subsection{\texorpdfstring{Fixed-point inverse}{Fixed-point inverse}}\label{subsubsec:simul-inversa}
% ---------------------------------------------------------------------
The inverse of $\mathbf{G}_c$ in fixed-point arithmetic is derived from the C++ code of the inversion module (Section~\ref{subsec:hw-modulo}). This code is compiled using \texttt{g++} along with the \texttt{ap\_fixed} header files from Vitis HLS 2026.1, with a separate compilation performed for each combination of format $\langle W, I\rangle$, quantization mode, overflow mode (Table~\ref{tab:modos-apfixed}), and size $M_c$. The program remains active throughout the simulation. Other calculations within the simulator, including the inverses used for BDAC initialization and LMMSE, utilize 64-bit floating-point arithmetic.

To verify this program, the same $\mathbf{G}_c$ matrices used in the simulator were inverted both by this program and by the Vitis HLS C-simulation. This process involved 30 matrices for each of the 35 mode combinations, with $M_c$ set to 16 and 32. In every instance, the resulting inverses were bit-identical. Consequently, the Symbol Error Rate (SER) is measured using the exact same inverse as that calculated by the hardware. The simulator also includes a Python-based model of the inversion, which has been verified bit-for-bit against the C++ code; however, the reported SER results are based on the C++ implementation.

% REMOVIDO (2026-11-05): subseção "Flow diagram". A Figura fig:fluxo_simulacao foi para a
%   abertura da seção; a frase da subseção ficou redundante com a abertura.



% =====================================================================
\section{\texorpdfstring{Hardware setup}{Hardware setup}}\label{sec:setup-hardware}
% =====================================================================
% =====================================================================
% SETUP DE HARDWARE (Metodologia) -- rascunho de 2026-11-04
% Ainda não incluído no METODOLOGIA.tex. Para incluir:
%   \section{Hardware Setup}\label{sec:setup-hardware}
%   % =====================================================================
% SETUP DE HARDWARE (Metodologia) -- rascunho de 2026-11-04
% Ainda não incluído no METODOLOGIA.tex. Para incluir:
%   \section{Hardware Setup}\label{sec:setup-hardware}
%   % =====================================================================
% SETUP DE HARDWARE (Metodologia) -- rascunho de 2026-11-04
% Ainda não incluído no METODOLOGIA.tex. Para incluir:
%   \section{Hardware Setup}\label{sec:setup-hardware}
%   \input{chapters/METODOLOGIA/setup-hardware.tex}
% Fontes do projeto (conferidas em 2026-11-04):
%   hardware/pacote_servidor_formatos/{LEIAME.md, vivado.tcl, potencia.py, rodar.sh}
%   hardware/pacote_servidor_2023/LEIAME.md
%   hardware/testes/{AMOSTRAGEM.md, LEIAME.md}
%   hardware/cholesky_hardware/{hls_config_sintese_24_12*.cfg, tb_potencia.cpp}
%   hardware/cholesky_hardware_diretivas*/LEIAME.md; acompanhamento_experimentos.csv
% Nomes das configurações: mapa LETRAS de software/testes/mapa_cenarios_v4.py
%   (antigos 1-8 = 1A,1E,1B,1F,1C,1G,1D,1H; 9-12 = 2A,2E,2B,2F).
% =====================================================================

% ---------------------------------------------------------------------
\subsection{\texorpdfstring{Tools and Target Device}{Tools and Target Device}}\label{subsec:hw-ferramentas}
% ---------------------------------------------------------------------
The inversion module was described in C++ and synthesized in Vitis HLS 2026.1. Vitis HLS generates the module's RTL. Then, the RTL was synthesized, placed, and routed in Vivado 2023.2. The target device is the AMD Alveo U50 board, with the XCU50-FSVH2104-2-E FPGA. This FPGA has 5,952 DSP blocks and 871,680 LUTs. All circuits use a clock frequency of 100 MHz, or a period of 10 ns.

In Vivado, the inversion module is synthesized and implemented alone, as the top-level design, without any other surrounding board logic. The flow performs synthesis, optimization, placement, physical optimization, and routing. The only timing constraint is the clock period, and the paths entering and exiting the module's ports are not constrained. Thus, the resource, \textcolor{red}{timing}, and power figures for this work are only for the module and come from the routed design in Vivado 2023.2.

Before the Vivado workflow, each version went through the Vitis HLS C synthesis. A version only went to Vivado if the HLS estimate fit within the FPGA: a maximum of 5,952 DSP blocks and a maximum of 90\% of the LUTs.

% ---------------------------------------------------------------------
\subsection{\texorpdfstring{Inversion Module}{Inversion Module}}\label{subsec:hw-modulo}
% ---------------------------------------------------------------------
The module receives the $M_c \times M_c$ matrix $\mathbf{G}_c$ and returns its inverse. It calculates the inverse using Cholesky decomposition (Section~\ref{sec:inversao}). Two sizes for $\mathbf{G}_c$, corresponding to the configurations described in Section~\ref{subsec:configuracoes-bs}, have been implemented: $16 \times 16$ and $32 \times 32$. The input, output, and internal values use the \texttt{ap\_fixed} data type, following the formats selected in Section~\ref{sec:formato}.

The module consists of six stages, each implemented as a C++ function (Figure~\ref{fig:modulo-inversao}). These stages are executed sequentially, with each stage passing its output to the next.
\begin{figure}[!ht]
\caption{Steps of the inversion module. Each block is one C++ function of the HLS code.}
\label{fig:modulo-inversao}
\vspace{4pt} % a classe zera o espaço entre a caption e a figura
\centering
\begin{tikzpicture}[
  font=\fontsize{7.5}{9}\selectfont,
  >={Stealth[length=2.1mm, width=1.7mm]},
  etapa/.style={draw=black!65, fill=white, align=center, text width=3.0cm,
                minimum height=1.1cm, inner sep=3pt},
  seta/.style={->, draw=black, line width=0.6pt},
  rot/.style={font=\fontsize{7}{8.4}\selectfont, inner sep=2pt},
]
% Duas linhas em U: 1 -> 2 -> 3 na linha de cima, desce, 4 -> 5 -> 6 na de baixo.
\newcommand{\fnome}[1]{{\fontsize{6.5}{7.8}\selectfont\ttfamily #1}}
\node[etapa] (e1) at (1.9,0)   {\fnome{LU\_decomposition}\\[2pt] $\mathbf{G}_c = \mathbf{L}\mathbf{U}$};
\node[etapa] (e2) at (7.3,0)   {\fnome{D\_extraction}\\[2pt] $\mathbf{D} = \mathrm{diag}(\mathbf{U})$};
\node[etapa] (e3) at (12.7,0)  {\fnome{D\_square\_root}\\[2pt] $\mathbf{D}^{1/2}$};
\node[etapa] (e4) at (12.7,-2.1) {\fnome{G\_decomposition}\\[2pt] $\mathbf{F} = \mathbf{L}\mathbf{D}^{1/2}$};
\node[etapa] (e5) at (7.3,-2.1)  {\fnome{forward\_substitution}\\[2pt] $\mathbf{F}^{-1}$};
\node[etapa] (e6) at (1.9,-2.1)  {\fnome{inversion}\\[2pt] $\mathbf{G}_c^{-1} = \mathbf{F}^{-H}\mathbf{F}^{-1}$};

\draw[seta] (e1) -- node[rot, above] {$\mathbf{U}$} (e2);
\draw[seta] (e2) -- node[rot, above] {$\mathbf{D}$} (e3);
\draw[seta] (e3) -- node[rot, right] {$\mathbf{D}^{1/2}$} (e4);
\draw[seta] (e4) -- node[rot, above] {$\mathbf{F}$} (e5);
\draw[seta] (e5) -- node[rot, above] {$\mathbf{F}^{-1}$} (e6);
% L vai direto da etapa 1 para a etapa 4, pelo espaço entre as linhas
\path let \p1=(e1.south), \p2=(e4.north), \n1={(\y1+\y2)/2} in
  coordinate (meio) at (0,\n1);
\draw[seta] (e1.south) -- (e1.south |- meio) -- node[rot, above, pos=0.5] {$\mathbf{L}$} ($(e4.north |- meio)+(-0.9,0)$) -- ($(e4.north)+(-0.9,0)$);

% entrada e saída
\draw[seta] (e1.west) ++(-1.0,0) node[rot, above right=0pt and -2pt] {$\mathbf{G}_c$} -- (e1.west);
\draw[seta] (e6.west) -- ++(-1.0,0) node[rot, above right=0pt and -2pt] {$\mathbf{G}_c^{-1}$};
\end{tikzpicture}
\vspace{4pt}
\legend{Source: the author.}
\end{figure}

Additions, multiplications, and divisions utilize \texttt{ap\_fixed} type operators. Complex multiplication and division operations are explicitly implemented in the code, with a specific type defined for each intermediate value; thus, the bit-width of every intermediate value is explicitly set within the code. The addition of two $N$-bit numbers requires $N+1$ bits, while their product requires $2N$ bits \cite[pp.~72--73]{kastner-hls}. When using \texttt{ap\_fixed} types, the result of an operation inherently possesses a sufficient number of bits; overflow only occurs when this result is assigned to a type with a smaller bit-width \cite[p.~74]{kastner-hls}. Consequently, intermediate values employ types with greater bit-widths than the $\langle W, I\rangle$ format (Table~\ref{tab:tipos-intermediarios}), with the result reverting to the specified format only at the output of each stage. Square root calculations are performed using the \texttt{hls::sqrt} function from Vitis HLS, operating directly within the defined format.

\begin{table}[!ht]
\caption{Types of the values in the inversion module, for the format $\langle W, I\rangle$ with $F = W - I$ fractional bits.}
\label{tab:tipos-intermediarios}
\vspace{4pt} % a classe zera o espaço entre a caption e a tabela
\centering
\footnotesize
\renewcommand{\arraystretch}{1.2}
\begin{tabular}{p{7.4cm}cc}
\toprule
\textbf{Value} & \textbf{Integer bits} & \textbf{Fractional bits} \\
\midrule
Input, output and output of each step & $I$ & $F$ \\
Square root (\texttt{hls::sqrt}) & $I$ & $F$ \\
Accumulated sums, numerator and divisor of the complex division, quotient & $2I$ & $2F$ \\
Product of two real numbers & $4I$ & $4F$ \\
Real or imaginary part of a complex product (sum of two products) & $4I + 1$ & $4F$ \\
\bottomrule
\end{tabular}
\vspace{4pt}
\legend{Source: the author. The intermediate values use the same quantization and overflow modes as the format; the square root uses AP\_TRN and AP\_WRAP.}
\end{table}
% ANTIGO (2026-11-05, versão em português, em vermelho):
% \begin{table}[!ht]
% \caption{\textcolor{red}{Tipos dos valores no módulo de inversão, para o formato $\langle W, I\rangle$ com $F = W - I$ bits fracionários.}}
% \label{tab:tipos-intermediarios}
% \centering
% \footnotesize
% \color{red}
% \renewcommand{\arraystretch}{1.2}
% \begin{tabular}{p{7.4cm}cc}
% \toprule
% \textbf{Valor} & \textbf{Bits inteiros} & \textbf{Bits fracionários} \\
% \midrule
% Entrada, saída e saída de cada etapa & $I$ & $F$ \\
% Raiz quadrada (\texttt{hls::sqrt}) & $I$ & $F$ \\
% Somas acumuladas, numerador e divisor da divisão complexa, quociente & $2I$ & $2F$ \\
% Produto de dois números reais & $4I$ & $4F$ \\
% Parte real ou imaginária de um produto complexo (soma de dois produtos) & $4I + 1$ & $4F$ \\
% \bottomrule
% \end{tabular}
% \vspace{4pt}
% \legend{Fonte: a autora. Os intermediários usam os mesmos modos de quantização e de estouro do formato; a raiz quadrada usa AP\_TRN e AP\_WRAP.}
% \end{table}


% ---------------------------------------------------------------------

% Fontes do projeto (conferidas em 2026-11-04):
%   hardware/pacote_servidor_formatos/{LEIAME.md, vivado.tcl, potencia.py, rodar.sh}
%   hardware/pacote_servidor_2023/LEIAME.md
%   hardware/testes/{AMOSTRAGEM.md, LEIAME.md}
%   hardware/cholesky_hardware/{hls_config_sintese_24_12*.cfg, tb_potencia.cpp}
%   hardware/cholesky_hardware_diretivas*/LEIAME.md; acompanhamento_experimentos.csv
% Nomes das configurações: mapa LETRAS de software/testes/mapa_cenarios_v4.py
%   (antigos 1-8 = 1A,1E,1B,1F,1C,1G,1D,1H; 9-12 = 2A,2E,2B,2F).
% =====================================================================

% ---------------------------------------------------------------------
\subsection{\texorpdfstring{Tools and Target Device}{Tools and Target Device}}\label{subsec:hw-ferramentas}
% ---------------------------------------------------------------------
The inversion module was described in C++ and synthesized in Vitis HLS 2026.1. Vitis HLS generates the module's RTL. Then, the RTL was synthesized, placed, and routed in Vivado 2023.2. The target device is the AMD Alveo U50 board, with the XCU50-FSVH2104-2-E FPGA. This FPGA has 5,952 DSP blocks and 871,680 LUTs. All circuits use a clock frequency of 100 MHz, or a period of 10 ns.

In Vivado, the inversion module is synthesized and implemented alone, as the top-level design, without any other surrounding board logic. The flow performs synthesis, optimization, placement, physical optimization, and routing. The only timing constraint is the clock period, and the paths entering and exiting the module's ports are not constrained. Thus, the resource, \textcolor{red}{timing}, and power figures for this work are only for the module and come from the routed design in Vivado 2023.2.

Before the Vivado workflow, each version went through the Vitis HLS C synthesis. A version only went to Vivado if the HLS estimate fit within the FPGA: a maximum of 5,952 DSP blocks and a maximum of 90\% of the LUTs.

% ---------------------------------------------------------------------
\subsection{\texorpdfstring{Inversion Module}{Inversion Module}}\label{subsec:hw-modulo}
% ---------------------------------------------------------------------
The module receives the $M_c \times M_c$ matrix $\mathbf{G}_c$ and returns its inverse. It calculates the inverse using Cholesky decomposition (Section~\ref{sec:inversao}). Two sizes for $\mathbf{G}_c$, corresponding to the configurations described in Section~\ref{subsec:configuracoes-bs}, have been implemented: $16 \times 16$ and $32 \times 32$. The input, output, and internal values use the \texttt{ap\_fixed} data type, following the formats selected in Section~\ref{sec:formato}.

The module consists of six stages, each implemented as a C++ function (Figure~\ref{fig:modulo-inversao}). These stages are executed sequentially, with each stage passing its output to the next.
\begin{figure}[!ht]
\caption{Steps of the inversion module. Each block is one C++ function of the HLS code.}
\label{fig:modulo-inversao}
\vspace{4pt} % a classe zera o espaço entre a caption e a figura
\centering
\begin{tikzpicture}[
  font=\fontsize{7.5}{9}\selectfont,
  >={Stealth[length=2.1mm, width=1.7mm]},
  etapa/.style={draw=black!65, fill=white, align=center, text width=3.0cm,
                minimum height=1.1cm, inner sep=3pt},
  seta/.style={->, draw=black, line width=0.6pt},
  rot/.style={font=\fontsize{7}{8.4}\selectfont, inner sep=2pt},
]
% Duas linhas em U: 1 -> 2 -> 3 na linha de cima, desce, 4 -> 5 -> 6 na de baixo.
\newcommand{\fnome}[1]{{\fontsize{6.5}{7.8}\selectfont\ttfamily #1}}
\node[etapa] (e1) at (1.9,0)   {\fnome{LU\_decomposition}\\[2pt] $\mathbf{G}_c = \mathbf{L}\mathbf{U}$};
\node[etapa] (e2) at (7.3,0)   {\fnome{D\_extraction}\\[2pt] $\mathbf{D} = \mathrm{diag}(\mathbf{U})$};
\node[etapa] (e3) at (12.7,0)  {\fnome{D\_square\_root}\\[2pt] $\mathbf{D}^{1/2}$};
\node[etapa] (e4) at (12.7,-2.1) {\fnome{G\_decomposition}\\[2pt] $\mathbf{F} = \mathbf{L}\mathbf{D}^{1/2}$};
\node[etapa] (e5) at (7.3,-2.1)  {\fnome{forward\_substitution}\\[2pt] $\mathbf{F}^{-1}$};
\node[etapa] (e6) at (1.9,-2.1)  {\fnome{inversion}\\[2pt] $\mathbf{G}_c^{-1} = \mathbf{F}^{-H}\mathbf{F}^{-1}$};

\draw[seta] (e1) -- node[rot, above] {$\mathbf{U}$} (e2);
\draw[seta] (e2) -- node[rot, above] {$\mathbf{D}$} (e3);
\draw[seta] (e3) -- node[rot, right] {$\mathbf{D}^{1/2}$} (e4);
\draw[seta] (e4) -- node[rot, above] {$\mathbf{F}$} (e5);
\draw[seta] (e5) -- node[rot, above] {$\mathbf{F}^{-1}$} (e6);
% L vai direto da etapa 1 para a etapa 4, pelo espaço entre as linhas
\path let \p1=(e1.south), \p2=(e4.north), \n1={(\y1+\y2)/2} in
  coordinate (meio) at (0,\n1);
\draw[seta] (e1.south) -- (e1.south |- meio) -- node[rot, above, pos=0.5] {$\mathbf{L}$} ($(e4.north |- meio)+(-0.9,0)$) -- ($(e4.north)+(-0.9,0)$);

% entrada e saída
\draw[seta] (e1.west) ++(-1.0,0) node[rot, above right=0pt and -2pt] {$\mathbf{G}_c$} -- (e1.west);
\draw[seta] (e6.west) -- ++(-1.0,0) node[rot, above right=0pt and -2pt] {$\mathbf{G}_c^{-1}$};
\end{tikzpicture}
\vspace{4pt}
\legend{Source: the author.}
\end{figure}

Additions, multiplications, and divisions utilize \texttt{ap\_fixed} type operators. Complex multiplication and division operations are explicitly implemented in the code, with a specific type defined for each intermediate value; thus, the bit-width of every intermediate value is explicitly set within the code. The addition of two $N$-bit numbers requires $N+1$ bits, while their product requires $2N$ bits \cite[pp.~72--73]{kastner-hls}. When using \texttt{ap\_fixed} types, the result of an operation inherently possesses a sufficient number of bits; overflow only occurs when this result is assigned to a type with a smaller bit-width \cite[p.~74]{kastner-hls}. Consequently, intermediate values employ types with greater bit-widths than the $\langle W, I\rangle$ format (Table~\ref{tab:tipos-intermediarios}), with the result reverting to the specified format only at the output of each stage. Square root calculations are performed using the \texttt{hls::sqrt} function from Vitis HLS, operating directly within the defined format.

\begin{table}[!ht]
\caption{Types of the values in the inversion module, for the format $\langle W, I\rangle$ with $F = W - I$ fractional bits.}
\label{tab:tipos-intermediarios}
\vspace{4pt} % a classe zera o espaço entre a caption e a tabela
\centering
\footnotesize
\renewcommand{\arraystretch}{1.2}
\begin{tabular}{p{7.4cm}cc}
\toprule
\textbf{Value} & \textbf{Integer bits} & \textbf{Fractional bits} \\
\midrule
Input, output and output of each step & $I$ & $F$ \\
Square root (\texttt{hls::sqrt}) & $I$ & $F$ \\
Accumulated sums, numerator and divisor of the complex division, quotient & $2I$ & $2F$ \\
Product of two real numbers & $4I$ & $4F$ \\
Real or imaginary part of a complex product (sum of two products) & $4I + 1$ & $4F$ \\
\bottomrule
\end{tabular}
\vspace{4pt}
\legend{Source: the author. The intermediate values use the same quantization and overflow modes as the format; the square root uses AP\_TRN and AP\_WRAP.}
\end{table}
% ANTIGO (2026-11-05, versão em português, em vermelho):
% \begin{table}[!ht]
% \caption{\textcolor{red}{Tipos dos valores no módulo de inversão, para o formato $\langle W, I\rangle$ com $F = W - I$ bits fracionários.}}
% \label{tab:tipos-intermediarios}
% \centering
% \footnotesize
% \color{red}
% \renewcommand{\arraystretch}{1.2}
% \begin{tabular}{p{7.4cm}cc}
% \toprule
% \textbf{Valor} & \textbf{Bits inteiros} & \textbf{Bits fracionários} \\
% \midrule
% Entrada, saída e saída de cada etapa & $I$ & $F$ \\
% Raiz quadrada (\texttt{hls::sqrt}) & $I$ & $F$ \\
% Somas acumuladas, numerador e divisor da divisão complexa, quociente & $2I$ & $2F$ \\
% Produto de dois números reais & $4I$ & $4F$ \\
% Parte real ou imaginária de um produto complexo (soma de dois produtos) & $4I + 1$ & $4F$ \\
% \bottomrule
% \end{tabular}
% \vspace{4pt}
% \legend{Fonte: a autora. Os intermediários usam os mesmos modos de quantização e de estouro do formato; a raiz quadrada usa AP\_TRN e AP\_WRAP.}
% \end{table}


% ---------------------------------------------------------------------

% Fontes do projeto (conferidas em 2026-11-04):
%   hardware/pacote_servidor_formatos/{LEIAME.md, vivado.tcl, potencia.py, rodar.sh}
%   hardware/pacote_servidor_2023/LEIAME.md
%   hardware/testes/{AMOSTRAGEM.md, LEIAME.md}
%   hardware/cholesky_hardware/{hls_config_sintese_24_12*.cfg, tb_potencia.cpp}
%   hardware/cholesky_hardware_diretivas*/LEIAME.md; acompanhamento_experimentos.csv
% Nomes das configurações: mapa LETRAS de software/testes/mapa_cenarios_v4.py
%   (antigos 1-8 = 1A,1E,1B,1F,1C,1G,1D,1H; 9-12 = 2A,2E,2B,2F).
% =====================================================================

% ---------------------------------------------------------------------
\subsection{\texorpdfstring{Tools and Target Device}{Tools and Target Device}}\label{subsec:hw-ferramentas}
% ---------------------------------------------------------------------
The inversion module was described in C++ and synthesized in Vitis HLS 2026.1. Vitis HLS generates the module's RTL. Then, the RTL was synthesized, placed, and routed in Vivado 2023.2. The target device is the AMD Alveo U50 board, with the XCU50-FSVH2104-2-E FPGA. This FPGA has 5,952 DSP blocks and 871,680 LUTs. All circuits use a clock frequency of 100 MHz, or a period of 10 ns.

In Vivado, the inversion module is synthesized and implemented alone, as the top-level design, without any other surrounding board logic. The flow performs synthesis, optimization, placement, physical optimization, and routing. The only timing constraint is the clock period, and the paths entering and exiting the module's ports are not constrained. Thus, the resource, \textcolor{red}{timing}, and power figures for this work are only for the module and come from the routed design in Vivado 2023.2.

Before the Vivado workflow, each version went through the Vitis HLS C synthesis. A version only went to Vivado if the HLS estimate fit within the FPGA: a maximum of 5,952 DSP blocks and a maximum of 90\% of the LUTs.

% ---------------------------------------------------------------------
\subsection{\texorpdfstring{Inversion Module}{Inversion Module}}\label{subsec:hw-modulo}
% ---------------------------------------------------------------------
The module receives the $M_c \times M_c$ matrix $\mathbf{G}_c$ and returns its inverse. It calculates the inverse using Cholesky decomposition (Section~\ref{sec:inversao}). Two sizes for $\mathbf{G}_c$, corresponding to the configurations described in Section~\ref{subsec:configuracoes-bs}, have been implemented: $16 \times 16$ and $32 \times 32$. The input, output, and internal values use the \texttt{ap\_fixed} data type, following the formats selected in Section~\ref{sec:formato}.

The module consists of six stages, each implemented as a C++ function (Figure~\ref{fig:modulo-inversao}). These stages are executed sequentially, with each stage passing its output to the next.
\begin{figure}[!ht]
\caption{Steps of the inversion module. Each block is one C++ function of the HLS code.}
\label{fig:modulo-inversao}
\vspace{4pt} % a classe zera o espaço entre a caption e a figura
\centering
\begin{tikzpicture}[
  font=\fontsize{7.5}{9}\selectfont,
  >={Stealth[length=2.1mm, width=1.7mm]},
  etapa/.style={draw=black!65, fill=white, align=center, text width=3.0cm,
                minimum height=1.1cm, inner sep=3pt},
  seta/.style={->, draw=black, line width=0.6pt},
  rot/.style={font=\fontsize{7}{8.4}\selectfont, inner sep=2pt},
]
% Duas linhas em U: 1 -> 2 -> 3 na linha de cima, desce, 4 -> 5 -> 6 na de baixo.
\newcommand{\fnome}[1]{{\fontsize{6.5}{7.8}\selectfont\ttfamily #1}}
\node[etapa] (e1) at (1.9,0)   {\fnome{LU\_decomposition}\\[2pt] $\mathbf{G}_c = \mathbf{L}\mathbf{U}$};
\node[etapa] (e2) at (7.3,0)   {\fnome{D\_extraction}\\[2pt] $\mathbf{D} = \mathrm{diag}(\mathbf{U})$};
\node[etapa] (e3) at (12.7,0)  {\fnome{D\_square\_root}\\[2pt] $\mathbf{D}^{1/2}$};
\node[etapa] (e4) at (12.7,-2.1) {\fnome{G\_decomposition}\\[2pt] $\mathbf{F} = \mathbf{L}\mathbf{D}^{1/2}$};
\node[etapa] (e5) at (7.3,-2.1)  {\fnome{forward\_substitution}\\[2pt] $\mathbf{F}^{-1}$};
\node[etapa] (e6) at (1.9,-2.1)  {\fnome{inversion}\\[2pt] $\mathbf{G}_c^{-1} = \mathbf{F}^{-H}\mathbf{F}^{-1}$};

\draw[seta] (e1) -- node[rot, above] {$\mathbf{U}$} (e2);
\draw[seta] (e2) -- node[rot, above] {$\mathbf{D}$} (e3);
\draw[seta] (e3) -- node[rot, right] {$\mathbf{D}^{1/2}$} (e4);
\draw[seta] (e4) -- node[rot, above] {$\mathbf{F}$} (e5);
\draw[seta] (e5) -- node[rot, above] {$\mathbf{F}^{-1}$} (e6);
% L vai direto da etapa 1 para a etapa 4, pelo espaço entre as linhas
\path let \p1=(e1.south), \p2=(e4.north), \n1={(\y1+\y2)/2} in
  coordinate (meio) at (0,\n1);
\draw[seta] (e1.south) -- (e1.south |- meio) -- node[rot, above, pos=0.5] {$\mathbf{L}$} ($(e4.north |- meio)+(-0.9,0)$) -- ($(e4.north)+(-0.9,0)$);

% entrada e saída
\draw[seta] (e1.west) ++(-1.0,0) node[rot, above right=0pt and -2pt] {$\mathbf{G}_c$} -- (e1.west);
\draw[seta] (e6.west) -- ++(-1.0,0) node[rot, above right=0pt and -2pt] {$\mathbf{G}_c^{-1}$};
\end{tikzpicture}
\vspace{4pt}
\legend{Source: the author.}
\end{figure}

Additions, multiplications, and divisions utilize \texttt{ap\_fixed} type operators. Complex multiplication and division operations are explicitly implemented in the code, with a specific type defined for each intermediate value; thus, the bit-width of every intermediate value is explicitly set within the code. The addition of two $N$-bit numbers requires $N+1$ bits, while their product requires $2N$ bits \cite[pp.~72--73]{kastner-hls}. When using \texttt{ap\_fixed} types, the result of an operation inherently possesses a sufficient number of bits; overflow only occurs when this result is assigned to a type with a smaller bit-width \cite[p.~74]{kastner-hls}. Consequently, intermediate values employ types with greater bit-widths than the $\langle W, I\rangle$ format (Table~\ref{tab:tipos-intermediarios}), with the result reverting to the specified format only at the output of each stage. Square root calculations are performed using the \texttt{hls::sqrt} function from Vitis HLS, operating directly within the defined format.

\begin{table}[!ht]
\caption{Types of the values in the inversion module, for the format $\langle W, I\rangle$ with $F = W - I$ fractional bits.}
\label{tab:tipos-intermediarios}
\vspace{4pt} % a classe zera o espaço entre a caption e a tabela
\centering
\footnotesize
\renewcommand{\arraystretch}{1.2}
\begin{tabular}{p{7.4cm}cc}
\toprule
\textbf{Value} & \textbf{Integer bits} & \textbf{Fractional bits} \\
\midrule
Input, output and output of each step & $I$ & $F$ \\
Square root (\texttt{hls::sqrt}) & $I$ & $F$ \\
Accumulated sums, numerator and divisor of the complex division, quotient & $2I$ & $2F$ \\
Product of two real numbers & $4I$ & $4F$ \\
Real or imaginary part of a complex product (sum of two products) & $4I + 1$ & $4F$ \\
\bottomrule
\end{tabular}
\vspace{4pt}
\legend{Source: the author. The intermediate values use the same quantization and overflow modes as the format; the square root uses AP\_TRN and AP\_WRAP.}
\end{table}
% ANTIGO (2026-11-05, versão em português, em vermelho):
% \begin{table}[!ht]
% \caption{\textcolor{red}{Tipos dos valores no módulo de inversão, para o formato $\langle W, I\rangle$ com $F = W - I$ bits fracionários.}}
% \label{tab:tipos-intermediarios}
% \centering
% \footnotesize
% \color{red}
% \renewcommand{\arraystretch}{1.2}
% \begin{tabular}{p{7.4cm}cc}
% \toprule
% \textbf{Valor} & \textbf{Bits inteiros} & \textbf{Bits fracionários} \\
% \midrule
% Entrada, saída e saída de cada etapa & $I$ & $F$ \\
% Raiz quadrada (\texttt{hls::sqrt}) & $I$ & $F$ \\
% Somas acumuladas, numerador e divisor da divisão complexa, quociente & $2I$ & $2F$ \\
% Produto de dois números reais & $4I$ & $4F$ \\
% Parte real ou imaginária de um produto complexo (soma de dois produtos) & $4I + 1$ & $4F$ \\
% \bottomrule
% \end{tabular}
% \vspace{4pt}
% \legend{Fonte: a autora. Os intermediários usam os mesmos modos de quantização e de estouro do formato; a raiz quadrada usa AP\_TRN e AP\_WRAP.}
% \end{table}


% ---------------------------------------------------------------------


% =====================================================================
\section{\texorpdfstring{\textcolor{red}{Configurations and test conditions}}{Configurations and test conditions}}\label{sec:configuracoes}
% =====================================================================
\subsection{\texorpdfstring{\textcolor{red}{Base station configurations}}{Base station configurations}}\label{subsec:configuracoes-bs}
The uplink of the massive MIMO system was investigated using two types of parameters. Hard constraints refer to the physical characteristics of the base station: $M = 128$ or 256 antennas, $C = 8$ or 16 clusters, and the size $M_c = M/C$ of the inversion module. Soft constraints depend on the UEs being served and the level of interference: $K = 8$ or 12 UEs, QPSK or 16-QAM modulation, and an IoTh ratio of 10 dB as a baseline case, or 20 dB representing an extreme, high-interference scenario. These configurations were adapted from \cite{zhao2025jsac}. Table~\ref{tab:paineis} presents these configurations.


Each configuration is identified by a label consisting of a number and a letter. The number denotes the hard constraints: 1 for $M = 128$ and $C = 8$; 2 for $M = 256$ and $C = 8$; and 3 for $M = 256$ and $C = 16$. The letter indicates the soft constraints and remains consistent across all three groups. \textcolor{red}{As configurações com o mesmo número formam um grupo: o grupo 1 tem as configurações 1A--1H, o grupo 2, as 2A--2H, e o grupo 3, as 3A--3H.}

\begin{table}[htbp]
\caption{\textcolor{red}{Configurações simuladas.}}
\label{tab:paineis}
\vspace{4pt} % a classe zera o espaço entre a caption e a tabela
\centering
\footnotesize
\renewcommand{\arraystretch}{1.2}
% sem espaço em volta das linhas do booktabs, para as bordas verticais não terem falhas
\setlength{\aboverulesep}{0pt}
\setlength{\belowrulesep}{0pt}
\begin{tabular}{|c|c|c|c|c|c|c|}
\toprule
\multirow{2}{*}{\textbf{Configuration}} & \multicolumn{3}{c|}{\textcolor{red}{\textbf{Hard constraints}}} & \multicolumn{3}{c|}{\textcolor{red}{\textbf{Soft constraints}}} \\
\cline{2-7}
 & {\boldmath $M$} & {\boldmath $C$} & {\boldmath $M_c$} & \textbf{IoTh (dB)} & \textbf{Modulation} & {\boldmath $K$} \\
-- & \texttt{M} & \texttt{C} & -- & \texttt{iot\_db} & \texttt{modulation} & \texttt{K} \\
\midrule
1A & \multirow{8}{*}{128} & \multirow{8}{*}{8} & \multirow{8}{*}{16} & \multirow{4}{*}{10} & \multirow{2}{*}{16-QAM} & 8 \\
1B & & & & & & 12 \\
\cline{6-7}
1C & & & & & \multirow{2}{*}{QPSK} & 8 \\
1D & & & & & & 12 \\
\cline{5-7}
1E & & & & \multirow{4}{*}{20} & \multirow{2}{*}{16-QAM} & 8 \\
1F & & & & & & 12 \\
\cline{6-7}
1G & & & & & \multirow{2}{*}{QPSK} & 8 \\
1H & & & & & & 12 \\
\midrule
2A & \multirow{8}{*}{256} & \multirow{8}{*}{8} & \multirow{8}{*}{32} & \multirow{4}{*}{10} & \multirow{2}{*}{16-QAM} & 8 \\
2B & & & & & & 12 \\
\cline{6-7}
2C & & & & & \multirow{2}{*}{QPSK} & 8 \\
2D & & & & & & 12 \\
\cline{5-7}
2E & & & & \multirow{4}{*}{20} & \multirow{2}{*}{16-QAM} & 8 \\
2F & & & & & & 12 \\
\cline{6-7}
2G & & & & & \multirow{2}{*}{QPSK} & 8 \\
2H & & & & & & 12 \\
\midrule
3A & \multirow{8}{*}{256} & \multirow{8}{*}{16} & \multirow{8}{*}{16} & \multirow{4}{*}{10} & \multirow{2}{*}{16-QAM} & 8 \\
3B & & & & & & 12 \\
\cline{6-7}
3C & & & & & \multirow{2}{*}{QPSK} & 8 \\
3D & & & & & & 12 \\
\cline{5-7}
3E & & & & \multirow{4}{*}{20} & \multirow{2}{*}{16-QAM} & 8 \\
3F & & & & & & 12 \\
\cline{6-7}
3G & & & & & \multirow{2}{*}{QPSK} & 8 \\
3H & & & & & & 12 \\
\bottomrule
\end{tabular}
\vspace{4pt}
\legend{\textcolor{red}{Fonte: configurações 1A, 1B, 1C, 1E, 2A e 3A adaptadas de \citesp{zhao2025jsac}; as outras, a autora. Em todas as configurações, $r = 8$. Nomes no código: campos de \texttt{Panel} (\texttt{dbp\_sim/params.py}).}}
% EN (versão anterior, 2026-11-05): Source: configurations 1A, 1B, 1C, 1E, 2A and 3A adapted from \citesp{zhao2025jsac}; the others, the author. In all configurations, $r = 8$. Names in the code: fields of \texttt{Panel} (\texttt{dbp\_sim/params.py}).
\end{table}

% AUTORA (9-results.tex:245) -- parágrafo retirado do texto (tom); o conteúdo ficou em comentário.
% Configurações 1, 9, 13, 3, 2, 5 = painéis (a), (b), (c), (e), (f), (g) da Fig. 8 de Zhao et al.
% O painel (d) não aparece no artigo. Na configuração 3 (painel e), r = 8 foi
% mantido; L577 liga r a K, então r = 12 também seria defensável.




\subsection{Test Conditions}\label{subsec:setup_condicoes}
% Criado em 2026-11-04 (reorganização pedida pela autora): junta o que define as condições
%   de teste. Origens: RESULTADOS/setup_simulacao.tex (subsubseção "Test Conditions", texto da
%   autora, e Tabela tab:condicoes_teste); METODOLOGIA/simulacao.tex (parágrafo do N);
%   METODOLOGIA/formato.tex (TDLA30-5 e TDLA30-10 com as mesmas curvas).

The simulated test conditions of Table~\ref{tab:n-maximo} define the carrier frequency and the number of noise samples in the Python simulator, as shown in Table~\ref{tab:condicoes_teste}. The simulation does not include a time axis, and the channel does not account for the Doppler effect, so only these two variables change between test conditions. A frequency of 3.5 GHz was used for FR1 conditions, within the n78 band (3300 to 3800 MHz, TDD), and 28 GHz was used for FR2 conditions, within the n257 band (26500 to 29500 MHz, TDD) \cite[Tabs.~5.2-1 and 5.2-2]{3gpp-ts38104}. In each realization, the channels at both frequencies use the same UE positions. % TODO: 3,5 GHz não tem fonte. Sugestão: citar uma referência que use 3,5 GHz como banda típica do 5G NR (ex.: livro 5G and Beyond), ou dizer só que é uma escolha deste trabalho.


\begin{table}[htbp]
\caption{Carrier frequency and number of noise samples in each test condition.}
\label{tab:condicoes_teste}
\vspace{4pt} % a classe zera o espaço entre a caption e a tabela
\centering
\footnotesize
\renewcommand{\arraystretch}{1.2}
\begin{tabular}{cccc}
\toprule
\textbf{Test condition} & \textbf{Frequency range} & \textbf{Carrier frequency (GHz)} & \textbf{Noise samples {\boldmath $N$}} \\
\texttt{M1\_Driver.} & -- & \texttt{M1\_Driver.} & \texttt{M5\_Noise.} \\
\texttt{test\_conditions} & & \texttt{test\_conditions} & \texttt{n\_noise\_samples} \\
\midrule
TDLA30-5 & FR1 & 3.5 & 192 \\
TDLA30-10 & FR1 & 3.5 & 192 \\
TDLA30-75 & FR2 & 28 & 192 \\
TDLA30-300 & FR2 & 28 & 168 \\
TDLC300-100 & FR1 & 3.5 & 132 \\
TDLB100-400 & FR1 & 3.5 & 24 \\
\bottomrule
\end{tabular}

\vspace{4pt}
\legend{Source: \citesp{3gpp-ts38104} and the author.}
\end{table}

This work employs a double-symbol DM-RS, located at symbols 2, 3, 10, and 11 of each slot \cite[Tab.~6.4.1.1.3-4]{3gpp-ts38211}, as it enables the differentiation of DM-RS signals among $K = 8$ or 12 UEs (Section~\ref{subsubsec:disponibilidade-amostras}). The value of $N$ for each condition is taken from Table~\ref{tab:n-maximo}, with an upper limit of 192, the value used in \cite{zhao2025jsac}. This $N$ assumes that the receiver waits for samples arriving after the data symbol (Section~\ref{subsec:pipeline}). The HST350 and HST500 conditions were not simulated, as they correspond to $N = 0$.

TDLA30-5 and TDLA30-10 share the same carrier frequency and the same $N$ value. Consequently, they exhibit identical SER curves, and the simulations cover 5 distinct test conditions. The only difference between them lies in the inversion cycle limit (Section~\ref{subsec:restricao-vazao}).

The test conditions are grouped into two mobility profiles based on UE speed (Table~\ref{tab:tempo-coerencia}). The low-mobility profile includes TDLA30-5, TDLA30-10, and TDLA30-75, with speeds up to 3.1 km/h. The high-mobility profile comprises TDLA30-300, TDLC300-100, and TDLB100-400, with speeds ranging from 11.6 to 123.3 km/h. The ITU-R uses 3 km/h for indoor and pedestrian users, and 30 and 120 km/h for users in vehicles \cite[Sec.~8.2 and Tab.~5]{itu-m2412}. In each profile, the fastest condition is the most demanding: it has the smallest $T_c$ and the smallest $N$. This condition is TDLA30-75 in the low-mobility profile and TDLB100-400 in the high-mobility profile.

% FONTE ITU: Rep. ITU-R M.2412-0, Sec. 8.2 (p. 16): Indoor Hotspot "stationary and pedestrian",
%   Dense Urban "pedestrian and vehicular", Rural "pedestrian, vehicular and high speed vehicular";
%   Tab. 5a (p. 20): indoor 3 km/h; Tab. 5b (p. 23): indoor 3 km/h, in-car 30 km/h;
%   Tab. 5c (p. 25): indoor 3 km/h, 120 km/h e 500 km/h. Arquivo: books/R-REP-M.2412-2017-PDF-E.pdf.


% =====================================================================
\section{\texorpdfstring{Inversion cycle limit}{Inversion cycle limit}}\label{sec:requisitos-tempo}
% =====================================================================
% ---------------------------------------------------------------------
\subsection{Inversões por bloco de coerência}\label{subsec:inversoes-por-tc}
% ---------------------------------------------------------------------
% CONTEÚDO  Modelo B: uma estimativa R_cc por PRB, com as N amostras do próprio PRB.
%           Por DU e por T_c: N_PRB inversões de R_cc e N_PRB inversões de G_c.
%           G_c = E_s H_c H_c^H + R_cc não depende da iteração l: inverte uma vez
%           por PRB e é reaproveitada nas T iterações (não multiplicar por T).
%           N_PRB = 275 é o pior caso do artigo; pela norma, 273 (FR1, 30 kHz,
%           100 MHz) e 264 (FR2, 120 kHz, 400 MHz).
% FONTE     paper.md:577 (N = 192 por realização); :504 (G_c); :614 (refazer por T_c);
%           :640 (N_PRB = 275, "worst case in 5G NR"); TS 38.104 Tab. 5.3.2-1 e 5.3.2-2.
%           Planilha: tabela "Número de operações por PRB".
% ESTADO    a escrever

% ---------------------------------------------------------------------
\subsection{Módulos de inversão}\label{subsec:modulos-inversao}
% ---------------------------------------------------------------------
% CONTEÚDO  Três módulos separados, por simplificação: (1) R_cc, (2) G_c,
%           (3) K x K, só na DU C. Módulos 1 e 3 ficam no caminho do PRB; o
%           módulo 2 fica fora se terminar antes da iteração 1.
%           Ligar à coluna "Onde roda" da Tabela tab:bcd-complexidade.
% FONTE     paper.md:136-142 (Eq. 7, inversa K x K); :504 (G_c); decisão da autora.
% ESTADO    a escrever

% ---------------------------------------------------------------------
\subsection{Coleta das amostras e processamento}\label{subsec:pipeline}
% ---------------------------------------------------------------------
% CONTEÚDO  Pipeline: junta as N amostras (t_coleta), processa (t_proc) e
%           equaliza os dados guardados. Remeter à Tabela tab:n-maximo para N.
%           t_coleta = min(T_c, ceil(N/48) slots): 2 ms (FR1) e 0,5 ms (FR2) nos
%           cenários lentos com N = 192; T_c nos rápidos.
% FONTE     slot: TS 38.211 Tab. 4.3.2-1; DM-RS: TS 38.211 Tab. 6.4.1.1.3-3 e -4;
%           planilha: tabela "Coleta de amostras e orçamento por inversão"
%           (corrigir o arredondamento: usar T_c exato em T_c/N_PRB).
% ESTADO    a escrever

% ---------------------------------------------------------------------
\subsection{Restrição de vazão}\label{subsec:restricao-vazao}
% ---------------------------------------------------------------------
% CONTEÚDO  t_proc <= T_c. Como os módulos são separados, cada um precisa de
%           II_inv <= T_c / N_PRB, sem depender de M_c nem de T.
%           Tabela por condição de teste: T_c/N_PRB em µs e em ciclos (clock pendente).
%           Valores (N_PRB = 275): 90,9091; 45,4545; 6,0606; 1,5152; 4,5455;
%           1,1364; 0,1947; 0,1363 µs.
%           É uma restrição de intervalo (vazão), não de latência.
% FONTE     T_c: Tabela tab:tempo-coerencia; N_PRB: paper.md:640;
%           planilha: tabela "Tempo para operações por PRB".
% ESTADO    a escrever

% ---------------------------------------------------------------------
\subsection{Restrição de latência}\label{subsec:restricao-latencia}
% ---------------------------------------------------------------------
% CONTEÚDO  t_coleta + t_proc <= D, com D só como fórmula (sem fonte; 4 ms do
%           eMBB [R7] é limite de fora). Se D for fixado, a restrição aperta nos
%           cenários lentos: II <= (D - t_coleta - L_PRB) / N_PRB.
%           O módulo de G_c tem uma condição de latência própria: terminar antes
%           de A e b chegarem para a iteração 1 (sem número).
%           Latência por salto entre DUs fica de fora.
% FONTE     decisão da autora (2026-10-12); [R7] para o valor de eMBB, se citado.
% ESTADO    a escrever

% ---------------------------------------------------------------------
\subsection{Escopo e pior caso}\label{subsec:escopo-tempo}
% ---------------------------------------------------------------------
% CONTEÚDO  Pior caso: condições HST, um limite que o NR não alcança.
%           Fora do escopo: escalonador trocando usuários por slot; latência por
%           salto. A simulação não tem eixo de tempo: as curvas de SER valem para
%           cada N (testes com N = 12, 48, 144 e 192 via --conditions).
%           HST: velocidade nominal da norma; pela fórmula, 2334 Hz a 3,5 GHz
%           daria ~720 km/h (nota em modelo-de-canal.tex).
% FONTE     TS 38.104 Sec. G.3; dbp_sim/m1_driver.py (sem eixo de tempo).
% ESTADO    a escrever


% =====================================================================
\section{\texorpdfstring{Selection of the fixed-point format}{Escolha do formato de ponto fixo}}\label{sec:formato}
% =====================================================================
% ---------------------------------------------------------------------
% 1.5  ESCOLHA DO FORMATO DE PONTO FIXO          -> nova seção da Metodologia
% ---------------------------------------------------------------------
% ORIGEM: TCC/Análise bits.md (resumo da autora), com as correções de 2026-11-04:
%   nomes 1A-3D no lugar das letras A-P; IoTh; "condição de teste" no lugar de
%   "cenário"; critério de escolha lexicográfico (não a razão perda/ocupação);
%   5 formatos finais (não 3); AP_TRN_ZERO corrigido; tabela de velocidades
%   retirada (já está em FUNDAMENTACAO/modelo-de-canal.tex).

\subsection{\texorpdfstring{SNR increase}{SNR increase}}

To select the $\langle W, I\rangle$ format (Section~\ref{subsubsec:ponto-fixo}), the BCD-MMSE algorithm was simulated using both fixed-point and 64-bit floating-point matrix inversion. Only the inversion of $\mathbf G_c$ differs between the two simulations; however, the SNR and SER are measured for the entire uplink. Each format was evaluated based on the SNR increase—that is, the additional SNR required by the uplink to achieve an $\mathrm{SER} = 10^{-2}$ compared to the floating-point case. This increase corresponds to the horizontal distance between the two SER curves, which is the standard method for comparing two schemes \cite[p.~55]{tse-viswanath}.

Two thresholds were employed. A 0.1~dB threshold represents a near-ideal format, where the SER curve is virtually identical to that of the floating-point implementation. A 3~dB threshold serves as the maximum acceptable increase. An increase of 3~dB is equivalent to doubling the power \cite[p.~73]{tse-viswanath}. Consequently, at this threshold, each UE would need to transmit with twice the energy per symbol to compensate for the error introduced by the matrix inversion (Section~\ref{subsec:snr}).



\subsection{\texorpdfstring{\textcolor{red}{Varredura inicial}}{Varredura inicial}}
\textcolor{red}{A primeira etapa variou a parte inteira e a parte fracionária separadamente, com a outra parte fixa (Tabela~\ref{tab:varredura-inicial}). Cada tamanho foi testado nos 35 modos de quantização e de estouro (Tabela~\ref{tab:modos-apfixed}), com o BCD-MMSE em 1 e em 4 iterações.}


\begin{table}[htbp]
\caption{\textcolor{red}{Tamanhos testados na varredura inicial.}}
\label{tab:varredura-inicial}
\centering
\footnotesize
\begin{tabular}{lll}
\toprule
Varredura & Bits inteiros ($I$) & Bits fracionários ($F$) \\
\midrule
Parte inteira & 6, 8, 10, 12, 14, 16, 20, 24, 32 & 32 \\
Parte fracionária & 14 & 4, 6, 8, 10, 12, 16, 20, 24, 32 \\
\bottomrule
\end{tabular}
\legend{\textcolor{red}{Fonte: a autora.}}
\end{table}
% NOTA: 9 + 9 = 18 tamanhos, mas <46,14> (I = 14, F = 32) aparece nas duas
%   varreduras: 17 tamanhos por varredura de configuração e condição.

\subsection{\texorpdfstring{\textcolor{red}{Configurações de teste para as matrizes}}{Configurações de teste para as matrizes}}
\textcolor{red}{Na varredura inicial, foram usadas as 16 configurações cuja curva em ponto flutuante atinge $\mathrm{SER} = 10^{-2}$ entre $-10$ e 10~dB: 1A--1H, 2A, 2B, 2E, 2F, 3A, 3B, 3E e 3F. Rodar as 16 nas 5 condições de teste daria 80 casos, cada um com 17 tamanhos e 35 modos. Por isso, foram escolhidas só 5 configurações.}
% CONFERIDO (2026-11-06): metrica_ponto_operacao.csv tem os 80 casos, pontos de operação de
%   -10,0 a 2,36 dB. 1C (configuração 5) no (192; 3,5 GHz) cruza no limite: SER = 1,00e-2 em
%   -10 dB. As QPSK com M = 256 (2C, 2D, 2G, 2H, 3C, 3D, 3G, 3H) cruzam abaixo de -10 dB
%   (curvas_qpsk_256, grade de -20 a 10 dB); a autora preferiu não citar isso por enquanto.
% ANTIGO (2026-11-06, texto reduzido; tirada a frase "As configurações QPSK dos grupos 2 e 3 foram incluídas depois"):
% \textcolor{red}{Na varredura inicial, havia 16 configurações: 1A--1H, 2A, 2B, 2E, 2F, 3A, 3B, 3E e 3F. As configurações QPSK dos grupos 2 e 3 foram incluídas depois. Simular as 16 configurações nas 5 condições de teste daria 80 varreduras, cada uma com 17 tamanhos e todos os modos. Por isso, a varredura inicial usou só 5 configurações.}
% ANTIGO (2026-11-06): o parágrafo começava em "Simular as 16 configurações ..." sem dizer quais eram.
% FONTE: software/testes/mapa_cenarios_v4.py, LETRAS (numeração antiga 1-16 = 1A-1H, 2A, 2E,
%   2B, 2F, 3A, 3E, 3B, 3F); a autora incluiu as QPSK dos grupos 2 e 3 por último.

\textcolor{red}{Os 80 casos foram ordenados pelo maior valor absoluto entre as partes reais e imaginárias de $\mathbf{G}_c$, pois o circuito guarda as duas partes em separado. Esse valor é chamado aqui de maior $|$parte$|$ de $\mathbf{G}_c$. Foram usadas as matrizes dos dois pontos de SNR em volta do ponto de operação, $\mathrm{SER} = 10^{-2}$ em ponto flutuante. Em cada caso, foi tomado o percentil 99, para que poucas matrizes raras não decidam a ordem.}
% ANTIGO (2026-11-06, texto reduzido; juntado com o parágrafo do percentil):
% \textcolor{red}{Para escolher essas configurações, os 80 casos (16 configurações nas 5 condições de teste) foram ordenados pelo tamanho dos elementos de $\mathbf{G}_c$. Em cada caso, o ponto de operação é a SNR em que a curva em ponto flutuante atinge $\mathrm{SER} = 10^{-2}$. Foram usadas as matrizes $\mathbf{G}_c$ dos dois pontos de SNR simulados em volta desse ponto. Os elementos de $\mathbf{G}_c$ são complexos, e o circuito guarda a parte real e a parte imaginária em separado. Por isso, de cada matriz foi tomado o maior valor absoluto entre as partes reais e imaginárias de todos os seus elementos. Esse valor é chamado aqui de maior $|$parte$|$ de $\mathbf{G}_c$.}

% (parágrafo do percentil juntado ao anterior)
% ANTIGO (2026-11-06, texto reduzido):
% \textcolor{red}{Cada caso tem 1600 ou 3200 matrizes. Para resumir esses valores em um número, foi usado o percentil 99: o valor que fica acima de 99\% dos valores do caso. Assim, só 1\% das matrizes tem um elemento maior que ele. O percentil foi usado no lugar do máximo para que poucas matrizes raras não decidam a ordem.}
% FONTE: software/testes/metrica_ponto_operacao.py (docstring, passos 1-4) e
%   resultados/4_varredura_formatos/metrica_ponto_operacao.csv (80 linhas, coluna gc_p99).
%   Ponto de operação: referência em ponto flutuante [1], BCD-MMSE com 1 iteração.
%   Matrizes por caso: 100 realizações x C DUs x 2 SNRs = 1600 (C = 8) ou 3200 (C = 16).
%   As configurações QPSK do grupo 1 (5 a 8) usam as matrizes das 16-QAM (1 a 4) nas suas
%   próprias SNRs: a modulação não muda G_c.
% "Partes em separado": hls_inverse / cholesky_ops.h usam ap_fixed para real e imaginária
%   (CONFERIR a forma de citar; a definição do formato está em FUNDAMENTACAO/inversao.tex).
% JUSTIFICATIVA do percentil 99 (no lugar do máximo): escrita por mim, não está registrada
%   pela autora. CONFERIR se é esse o motivo.

\textcolor{red}{As configurações 1G, 1H, 2E e 2F ocupam as 17 primeiras posições. Elas têm os maiores elementos e, por isso, precisam de mais bits inteiros. A configuração 1A tem os menores elementos e foi usada como caso base (Tabela~\ref{tab:gc-ponto-operacao}).}
% ANTIGO (2026-11-06, texto reduzido):
% \textcolor{red}{Quatro delas são as configurações em que a matriz $\mathbf{G}_c$, no ponto de operação $\mathrm{SER} = 10^{-2}$, tem os elementos de maior valor absoluto: 2E, 2F, 1G e 1H. Elementos maiores precisam de mais bits inteiros para não saturar a entrada da inversão. Assim, essas são as configurações mais exigentes. A quinta é a configuração 1A, que tem os elementos de menor valor absoluto. Ela foi usada como caso base, para comparação. A Tabela~\ref{tab:gc-ponto-operacao} mostra esses valores.}

% (parágrafo das posições juntado ao anterior)
% ANTIGO (2026-11-06, texto reduzido):
% \textcolor{red}{As quatro configurações exigentes aparecem repetidas no topo da ordenação. As configurações 1G e 1H ocupam as 10 primeiras posições, com todas as suas condições de teste. As configurações 2E e 2F ocupam as posições 11 a 17. Já a configuração 1A ocupa a última posição, e as suas outras condições ficam entre as posições 75 e 78.}
% FONTE: metrica_ponto_operacao.csv, coluna posicao. 1G = 6, 1H = 8: posições 1-10.
%   2E = 10, 2F = 12: posições 11-17 (as condições restantes em 20, 27 e 29; N = 24 nas
%   duas últimas). 1A = 1: posições 75-78 e 80 (79 é 1B com N = 24).
%   Logo depois de 2E e 2F vêm 3F (18) e 3E (19), com valores parecidos (948 e 944).

\begin{table}[htbp]
\caption{\textcolor{red}{SNR no ponto de operação e maior elemento de $\mathbf{G}_c$ nas configurações da varredura inicial.}}
\label{tab:gc-ponto-operacao}
\centering
\footnotesize
\begin{tabular}{ccccc}
\toprule
 & \multicolumn{2}{c}{SNR no ponto de operação (dB)} & \multicolumn{2}{c}{Maior $|$parte$|$ de $\mathbf{G}_c$ (percentil 99)} \\
\cmidrule(lr){2-3}\cmidrule(lr){4-5}
Configuração & $N = 192$; 3,5 GHz & $N = 168$; 28 GHz & $N = 192$; 3,5 GHz & $N = 168$; 28 GHz \\
\midrule
1A & $-2{,}6$ & $-2{,}4$ & 63   & 66   \\
\textcolor{red}{1G} & $-9{,}3$ & $-9{,}1$ & 2266 & 2383 \\
\textcolor{red}{1H} & $-8{,}7$ & $-8{,}6$ & 2470 & 2385 \\
\textcolor{red}{2E} & $-5{,}8$ & $-5{,}6$ & 962  & 999  \\
\textcolor{red}{2F} & $-5{,}6$ & $-5{,}5$ & 983  & 987  \\
\bottomrule
\end{tabular}
\legend{\textcolor{red}{Fonte: a autora.}}
% ANTIGO (legenda): "Parte" é a parte real ou imaginária de cada elemento. -> foi para o texto.
\end{table}
% CORREÇÃO (2026-11-04): as linhas estavam com os nomes trocados (2E/2F com os valores de
%   1G/1H e vice-versa). Pelo metrica_ponto_operacao.csv: config 6 (1G) = -9,29/-9,13 dB e
%   2265,6/2382,9; config 8 (1H) = -8,74/-8,59 e 2470,0/2385,2; config 10 (2E) = -5,80/-5,60
%   e 961,5/999,0; config 12 (2F) = -5,61/-5,49 e 983,4/986,5. As QPSK (1G, 1H) operam em
%   SNR mais baixa e por isso têm G_c maior. Ordem antiga das linhas: 1A, 2E, 2F, 1G, 1H.
% CORREÇÃO antiga de nomes (mapa LETRAS de software/testes/mapa_cenarios_v4.py):
%   A = config 1 -> 1A; M = 10 -> 2E; N = 12 -> 2F; P = 6 -> 1G; O = 8 -> 1H.
%   No Análise bits.md, a primeira tabela troca o K de O e P (O com K = 8, P com
%   K = 12). Pelos números, 6 (1G) tem K = 8 e 8 (1H) tem K = 12; a tabela do
%   grupo 16x16, mais abaixo no mesmo arquivo, está certa. As linhas de 1G e 1H
%   acima seguem os números (P = 1G, O = 1H). CONFERIR os valores no CSV de origem.
% A CONFERIR: a tabela mostra só as 5 escolhidas. Para sustentar "as 4 maiores",
%   vale mostrar (ou citar em nota) o percentil 99 das 16 configurações.
% Ressalva (memória do projeto): o máximo absoluto de |G_c| chegou a 3960 em 1G/1H
%   a -10 dB, contra o limite de 4096 de I = 13; por isso I = 12 satura.

% SEPARAÇÃO (2026-11-04): a subseção dos modos de quantização e de estouro (texto e
%   Tabela tab:modos-apfixed) foi para FUNDAMENTACAO/inversao.tex.

\subsection{\texorpdfstring{\textcolor{red}{Síntese e seleção dos formatos para a validação}}{Síntese e seleção dos formatos para a validação}}
\textcolor{red}{Os tamanhos e modos da varredura inicial também foram sintetizados no Vitis HLS, para medir a ocupação de recursos (LUT, DSP, FF e BRAM). A ocupação de cada circuito é a do recurso mais usado no FPGA. Para cada parte da palavra, foram escolhidos os tamanhos com menor ocupação em duas classes: aumento de SNR de até 3~dB e de até 0,1~dB. Combinando os tamanhos escolhidos para a parte inteira e para a fracionária, ficaram 10 formatos para a etapa de validação.}

\textcolor{red}{A escolha foi feita para cada par (configuração, condição de teste), separada por limite de aumento de SNR e por número de iterações. Quando dois tamanhos empatam na ocupação, fica o que se repete em mais pares. O formato de um par junta a parte inteira e a parte fracionária escolhidas para ele. Os 10 formatos são todos os que apareceram nessa escolha. A Tabela~\ref{tab:formatos-origem} mostra em quais configurações cada formato foi escolhido.}

\begin{table}[htbp]
\caption{\textcolor{red}{Configurações em que cada formato foi escolhido na varredura inicial.}}
\label{tab:formatos-origem}
\centering
\footnotesize
\begin{tabular}{ccccc}
\toprule
 & \multicolumn{2}{c}{1 iteração} & \multicolumn{2}{c}{4 iterações} \\
\cmidrule(lr){2-3}\cmidrule(lr){4-5}
Formato & Aumento até 3~dB & Aumento até 0,1~dB & Aumento até 3~dB & Aumento até 0,1~dB \\
\midrule
$\langle 14,6\rangle$  & 1A         & --                       & --             & --                       \\
$\langle 16,6\rangle$  & --         & --                       & 1A             & --                       \\
$\langle 18,8\rangle$  & --         & 1A                       & --             & --                       \\
$\langle 20,8\rangle$  & --         & --                       & --             & 1A                       \\
$\langle 22,12\rangle$ & 1H, 2E, 2F & --                       & --             & --                       \\
$\langle 24,12\rangle$ & 1G         & 2E, 2F                   & 1G, 1H, 2E, 2F & --                       \\
$\langle 26,12\rangle$ & --         & 1G$^{a}$, 1H$^{a}$       & --             & 2E, 2F                   \\
$\langle 28,12\rangle$ & --         & --                       & --             & 1G$^{a}$, 1H$^{a}$       \\
$\langle 28,14\rangle$ & --         & 1G$^{b}$, 1H$^{b}$       & --             & --                       \\
$\langle 30,14\rangle$ & --         & --                       & --             & 1G$^{b}$, 1H$^{b}$       \\
\bottomrule
\end{tabular}
\legend{\textcolor{red}{Fonte: a autora. Sem marca: nas 5 condições de teste. $^{a}$Só na TDLB100-400. $^{b}$Nas outras 4 condições.}}
\end{table}
% FONTE: resultados/4_varredura_formatos/v2_casos_extremos/restricoes_varredura_4_detalhe.csv
%   (restricoes_v2.py). Cada linha do CSV é um lado (I ou F) de um par; o formato do par
%   é <I + F, I>. Nomes: 1 = 1A, 6 = 1G, 8 = 1H, 10 = 2E, 12 = 2F. TDLA30-10 = TDLA30-5.
% DESEMPATE (restricoes_v2.py): 1) a opção (tamanho e modo) que empata em mais células;
%   2) no lado de F, AP_RND_ZERO; 3) menor perda; 4) ordem dos modos. O texto cita só o 1.
% A "repetição" só desempata. Nenhum formato que apareceu foi descartado: os 10 são a
%   união das escolhas.

\textcolor{red}{Os formatos menores vêm do caso base, a configuração 1A. Os formatos com 12 e 14 bits inteiros vêm das configurações mais exigentes. Os formatos com mais bits fracionários vêm do limite de 0,1~dB e de 4 iterações.}

% CORREÇÃO: o Análise bits.md fala em minimizar "aumento de SNR / % de ocupação".
%   O critério aplicado é em ordem: primeiro o limite de perda, depois a menor
%   ocupação. Não é uma razão.
% A CONFERIR: como os pares (I, F) foram combinados nos 10 formatos.

\subsection{\texorpdfstring{\textcolor{red}{Validação}}{Validação}}
\textcolor{red}{A varredura inicial fixou uma parte da palavra enquanto variava a outra, e usou só 5 configurações. As duas coisas podem mudar o recurso mais usado e o melhor modo. Por isso, os 10 formatos foram simulados de novo, com os 35 modos, nas \textcolor{red}{24} configurações e nas 5 condições de teste, com o BCD-MMSE em 1 iteração.}
% ANTIGO: "nas 16 configurações". As 8 configurações QPSK com M = 256 (2C, 2D, 2G, 2H, 3C,
%   3D, 3G, 3H) entraram na [4] v5, completa (1750 rodadas por configuração = 10 x 35 x 5).
% FONTE: software/testes/LEIAME.md, varredura [4] v4 (28.000 rodadas = 10 formatos x
%   35 modos x 16 configurações x 5 condições), só nas SNRs em volta de SER = 1e-2.

\textcolor{red}{O circuito tem um modo só. Por isso, cada formato recebeu um modo fixo \textcolor{red}{em cada grupo de configurações (Seção~\ref{sec:configuracoes}). Todas as configurações de um grupo usam o mesmo circuito.} O modo foi escolhido pelos seguintes critérios, nesta ordem:}
% ANTIGO: "cada formato recebeu um modo fixo para os circuitos 16 x 16 (configurações 1A--1H
%   e 3A--3H) e outro para os circuitos 32 x 32 (configurações 2A--2H)."
% MOTIVO: os mapas por grupo (mapa_grupos_v4.py) escolhem o modo só com as configurações do
%   grupo (modo_fixo(p, occ_k, ks, ...) com ks = configurações do grupo). Os grupos 1 e 3
%   têm o mesmo circuito 16 x 16, mas podem receber modos diferentes. RESOLVIDO (2026-11-04):
%   a autora decidiu fazer a análise por grupo (modo e formatos escolhidos por grupo).
\begin{enumerate}
\item \textcolor{red}{o modo que deixa mais pares (configuração, condição de teste) com aumento de SNR de até 3~dB. Um par que não atinge $\mathrm{SER} = 10^{-2}$ conta como fora do limite;}
\item \textcolor{red}{a menor ocupação, pelo recurso mais usado na síntese;}
\item \textcolor{red}{a quantização \texttt{AP\_RND\_ZERO};}
\item \textcolor{red}{o menor aumento de SNR no pior par;}
\item \textcolor{red}{a ordem dos modos na lista.}
\end{enumerate}
% FONTE: mapa_cenarios_v4.py, modo_fixo (linhas 149-176): a chave é (pares dentro do
%   limite, ocupação, AP_RND_ZERO, pior perda, ordem dos modos). RESOLVIDO: AP_RND_ZERO
%   vem antes da pior perda (item acrescentado em vermelho).

\subsection{\texorpdfstring{\textcolor{red}{Mapas de aumento de SNR}}{Mapas de aumento de SNR}}\label{subsec:mapas-perda}
\textcolor{red}{O resultado da validação é mostrado em mapas de aumento de SNR, um para cada grupo de configurações. Cada linha do mapa é um formato, com o modo fixo do grupo. Cada coluna é uma configuração. A célula mostra o maior aumento de SNR entre as condições de teste. Os símbolos mostram as condições com aumento de SNR acima de 0,1~dB (vazados) e acima de 3~dB (cheios). Um x marca a configuração em que alguma condição não atinge $\mathrm{SER} = 10^{-2}$. Os mapas estão no capítulo de resultados.}
% FIGURAS: resultados/4_varredura_formatos/v4_formatos_juntos_16_configs/figuras/por_grupo/
%   mapa_perda_modo_fixo_ate_3dB_iter1_grupo{1,2,3}_modo_por_grupo.pdf (mapa_grupos_v4.py).
%   Grupos 2 e 3 já com as QPSK da v5. Escala de cor logarítmica, de 0,01 a 3 dB.

\subsection{\texorpdfstring{\textcolor{red}{Formatos escolhidos}}{Formatos escolhidos}}
\textcolor{red}{Os formatos foram escolhidos em cada grupo de configurações, com o modo fixo do grupo. Em cada grupo, foram escolhidos três formatos: o primeiro em que todas as configurações e condições atingem $\mathrm{SER} = 10^{-2}$; o primeiro em que todas atingem esse ponto com aumento de SNR de até 3~dB; e o primeiro em que todas atingem com aumento de até 0,1~dB. Os formatos são percorridos do menor para o maior $W$. A Tabela~\ref{tab:formatos-escolhidos} mostra os formatos escolhidos. Esses formatos foram implementados no Vivado, para medir a potência e as arquiteturas de cada condição de teste.}

\begin{table}[htbp]
\caption{\textcolor{red}{Formatos escolhidos em cada grupo de configurações.}}
\label{tab:formatos-escolhidos}
\centering
\footnotesize
\setlength{\tabcolsep}{3pt}
\begin{tabular}{>{\raggedright\arraybackslash}p{2.1cm}ccc}
\toprule
Critério & Grupo 1 & Grupo 2 & Grupo 3 \\
\midrule
Atinge $\mathrm{SER} = 10^{-2}$ & \begin{tabular}[c]{@{}c@{}}$\langle 22,12\rangle$\\ \texttt{AP\_RND\_ZERO}/\texttt{AP\_WRAP}\end{tabular} & \begin{tabular}[c]{@{}c@{}}$\langle 24,12\rangle$\\ \texttt{AP\_RND\_ZERO}/\texttt{AP\_WRAP}\end{tabular} & \begin{tabular}[c]{@{}c@{}}$\langle 24,12\rangle$\\ \texttt{AP\_RND\_ZERO}/\texttt{AP\_WRAP}\end{tabular} \\
\addlinespace
Aumento de até 3~dB & \begin{tabular}[c]{@{}c@{}}$\langle 24,12\rangle$\\ \texttt{AP\_RND\_ZERO}/\texttt{AP\_WRAP}\end{tabular} & \begin{tabular}[c]{@{}c@{}}$\langle 28,14\rangle$\\ \texttt{AP\_RND\_ZERO}/\texttt{AP\_WRAP}\end{tabular} & \begin{tabular}[c]{@{}c@{}}$\langle 24,12\rangle$\\ \texttt{AP\_RND\_ZERO}/\texttt{AP\_WRAP}\end{tabular} \\
\addlinespace
Aumento de até 0,1~dB & \begin{tabular}[c]{@{}c@{}}$\langle 28,14\rangle$\\ \texttt{AP\_RND\_ZERO}/\texttt{AP\_WRAP}\end{tabular} & \begin{tabular}[c]{@{}c@{}}$\langle 30,14\rangle$\\ \texttt{AP\_TRN}/\texttt{AP\_SAT}\end{tabular} & \begin{tabular}[c]{@{}c@{}}$\langle 30,14\rangle$\\ \texttt{AP\_TRN}/\texttt{AP\_SAT}\end{tabular} \\
\bottomrule
\end{tabular}
\legend{\textcolor{red}{Fonte: a autora. Abaixo do formato, o modo de quantização e o de estouro.}}
\end{table}
% FONTE: mesmo cálculo de mapa_grupos_v4.py (modo_fixo com as configurações do grupo, limite
%   3 dB, 1 iteração, v4 + v5). Script de conferência no scratchpad da sessão 981d36d9
%   (formatos_por_grupo.py). Pior perda no grupo (todas as configurações e condições):
%   G1: <22,12> 3,33; <24,12> 0,45; <26,12> 0,16; <28,12> 0,17; <28,14> 0,02; <30,14> 0,01.
%   G2: <22,12> x; <24,12> 4,12; <26,12> 3,76; <28,12> 3,67; <28,14> 0,34; <30,14> 0,04.
%   G3: <22,12> x; <24,12> 2,55; <26,12> 2,45; <28,12> 2,46; <28,14> 0,13; <30,14> 0,03.
%   Formatos de 6 e 8 bits inteiros têm x em todos os grupos.
% "Grupo 1/2/3": os nomes G1-G3 eram só internos (memória tcc-nomes-configuracoes). Aqui o
%   grupo é o número da configuração; CONFERIR se a autora quer "Grupo 1" na tese ou outro
%   rótulo (ex.: "1A--1H").
% <26,12>: com a análise por grupo, não é escolhido em nenhum grupo.
% ANTIGO (lista por todas as configurações juntas, antes da análise por grupo):
% Da validação, foram escolhidos os formatos para a implementação no Vivado, a medição de potência e as arquiteturas de cada condição de teste:
% \begin{itemize}
% \item $\langle 22,12\rangle$: o primeiro formato em que todas as configurações e condições atingem $\mathrm{SER} = 10^{-2}$;
% \item $\langle 24,12\rangle$: o primeiro formato em que todas atingem $\mathrm{SER} = 10^{-2}$ com perda de até 3~dB;
% \item $\langle 28,14\rangle$: o primeiro formato em que todas atingem $\mathrm{SER} = 10^{-2}$ com perda de até 0,1~dB;
% \item $\langle 26,12\rangle$: % PENDENTE: critério da escolha (decidido em 2026-10-31).
% \item $\langle 30,14\rangle$ em \texttt{AP\_TRN}/\texttt{AP\_SAT}: % PENDENTE: critério.
%       % ATENÇÃO (1 iteração): a única justificativa registrada era atender T = 4
%       % (0,04 dB na [5]). Com T = 1, <28,14> já fica até 0,1 dB. Na primeira versão,
%       % decidir se <30,14> entra e por quê; na versão com 4 iterações (se houver
%       % tempo), essa é a justificativa natural.
% \end{itemize}
% Os quatro primeiros usam \texttt{AP\_RND\_ZERO}/\texttt{AP\_WRAP}.
% % FONTE: software/testes/LEIAME.md, rodada [5] modo_fixo_3dB; mapa iter1 (decisão de
% %   2026-11-04: o TCC usa só 1 iteração). "Primeiro" = menor W na ordem dos 10 formatos.
%


% =====================================================================
\section{\texorpdfstring{\textcolor{red}{Hardware architectures and tests}}{Hardware architectures and tests}}\label{sec:testes-hardware}
% =====================================================================
% Separado de setup-hardware.tex em 2026-11-04: as ferramentas e o módulo ficam antes
%   (ferramentas); as arquiteturas, as matrizes de teste e a potência vêm depois do formato.
\subsection{\texorpdfstring{\textcolor{red}{Hardware Architectures}}{Hardware Architectures}}\label{subsec:hw-arquiteturas}
% ---------------------------------------------------------------------
\textcolor{red}{Cada condição de teste define um intervalo máximo entre inversões, $T_c/N_{\text{PRB}}$ (Seção~\ref{subsec:restricao-vazao}). Com o clock de 100~MHz, esse intervalo é dado em ciclos. Foi feita uma arquitetura para cada condição de teste, com diretivas do HLS que reduzem o intervalo até ele atender esse limite. Também foi implementada uma versão sem otimizações de ciclos, como base de comparação. Essa versão não tem diretivas. Além disso, o Vitis HLS aplica \textit{pipeline} sozinho aos laços internos. Nessa versão, esse \textit{pipeline} automático foi desligado, com a opção \texttt{syn.compile.pipeline\_loops=0}. Assim, as iterações de cada laço rodam uma depois da outra. A Tabela~\ref{tab:hw-arquiteturas} mostra as arquiteturas.}
% ANTIGO (2026-11-06): ... como base de comparação. A Tabela~\ref{tab:hw-arquiteturas} mostra as arquiteturas.
% FONTE: hardware/cholesky_hardware/hls_config_sintese_24_12.cfg, linhas 9-10 e 16 ("pipeline_loops=0
%   desliga o pipeline automático que o Vitis aplica sozinho nos laços internos (padrão 64). Sem
%   isso a síntese já sai com 12 laços pipelinados"). TODO FONTE: guia do Vitis HLS (UG1399),
%   opção syn.compile.pipeline_loops, se a banca pedir referência do comportamento padrão.
% EN (versão anterior, 2026-11-05): Each test condition sets a maximum interval between inversions, $T_c/N_{\text{PRB}}$ (Section~\ref{subsec:restricao-vazao}). At 100~MHz, this interval is given in clock cycles. One architecture was designed for each test condition, with HLS directives that reduce the interval until it meets this limit. A version without cycle optimizations was also implemented, as a baseline. Table~\ref{tab:hw-arquiteturas} shows the architectures.

\begin{table}[htbp]
\caption{\textcolor{red}{Arquiteturas de hardware do módulo de inversão.}}
\label{tab:hw-arquiteturas}
\vspace{4pt}
\centering
\footnotesize
\color{red}% conteúdo da tabela traduzido para o português (2026-11-05); era em inglês
\renewcommand{\arraystretch}{1.2}
\begin{tabular}{lp{2.2cm}p{1.6cm}p{6.2cm}}
\toprule
\textbf{Arquitetura} & \textbf{Condição de teste alvo} & \textbf{Limite de ciclos} & \textbf{Principal otimização} \\
\midrule
Sem otimização & -- & -- & Sem diretivas; \textit{pipeline} automático de laços desligado \\
9090 & TDLA30-5 & 9.090 & Laços com \textit{pipeline} na decomposição, na substituição direta e na inversão \\
4545 & TDLA30-10 & 4.545 & A mesma, mais \textit{dataflow} entre as etapas \\
606 & TDLA30-75 & 606 & Lote de 12 matrizes em 6 vias, com \textit{dataflow} \\
455 & TDLC300-100 & 454 & Lote de 16 matrizes em 8 vias \\
152 & TDLA30-300 & 151 & Arranjo sistólico, lote de 32 matrizes em 2 vias \\
114 & TDLB100-400 & 113 & Arranjo sistólico, lote de 32 matrizes em 4 vias \\
\bottomrule
\end{tabular}
\vspace{4pt}
\legend{\textcolor{red}{Fonte: a autora. Limite de ciclos da Tabela~\ref{tab:limite-ciclos}.}}
% ANTIGO (2026-11-05): legenda "Limite de ciclos: T_c/N_PRB, com N_PRB = 275, a 100 MHz"; limites 152, 455 e 114
%   estavam arredondados para cima. Os nomes das arquiteturas (152, 455, 114) não mudaram.
% EN (versão anterior, 2026-11-05): Source: the author. Cycle limit: $T_c/N_{\text{PRB}}$ with $N_{\text{PRB}} = 275$, at 100~MHz.
\end{table}
% FONTE: acompanhamento_experimentos.csv (coluna "descrição da otimização aplicada");
%   LEIAME.md de cada pasta cholesky_hardware_diretivas*; pipeline_loops=0 em
%   hls_config_sintese_24_12*.cfg. Limites: memória tcc-constraints-tempo-estado.
% A CONFERIR: descrição da 4545 ("alterações 0 e 1 (dataflow)") e da 9090 ("alterações
%   1 a 3, 5 e 6") -- detalhar a partir dos LEIAME se a banca precisar.
% Os ciclos obtidos (7.594, 3.020, 504, 402, 128, 64 em <24,12>) são RESULTADO: vão para
%   os Resultados, com a origem (log xsim do servidor), e mudam com o formato.

\textcolor{red}{As arquiteturas otimizadas foram implementadas só em $16 \times 16$. Em $32 \times 32$, só a versão sem otimizações foi implementada no Vivado.}
% EN (versão anterior, 2026-11-05): The optimized architectures were implemented only in $16 \times 16$. In $32 \times 32$, only the version without optimizations was implemented in Vivado.
% FONTE: acompanhamento_experimentos.csv, linha "otimizações de ciclos (32 x 32)":
%   só síntese C local; nenhuma atende o alvo de ciclos; 114 não cabe. Isso é resultado:
%   a frase aqui só diz o que foi implementado; o porquê vai para os Resultados.

\textcolor{red}{Todas as arquiteturas calculam a mesma inversa que a versão sem otimizações. Isso foi conferido bit a bit, com o código C de cada versão, nos formatos escolhidos e nos dois tamanhos de matriz.}
% EN (versão anterior, 2026-11-05): Every architecture computes the same inverse as the version without optimizations. This was checked bit by bit, with the C code of each version, in the chosen formats and in both matrix sizes.
% FONTE: comparar_com_original.py (LEIAME de cada versão); conferir_formatos_novos.sh
%   para <26,12> e <30,14> (iguais nas 6 arquiteturas, 16x16 e 32x32).

% =====================================================================
% Subseções das otimizações (2026-11-06). Uma por tipo de otimização, com figura,
%   diretiva e trecho de código. Fontes: Kastner (books/Parallel Programming for
%   FPGAs.pdf) e o código de cada pasta hardware/cholesky_hardware_diretivas*.
% As figuras são esquemáticas: nenhum número medido aqui (os ciclos medidos são
%   resultado e vão para os Resultados). Os trechos de código são simplificados.
% =====================================================================
\subsubsection{\texorpdfstring{\textcolor{red}{Pipeline e tempo de enchimento}}{Pipeline e tempo de enchimento}}\label{subsubsec:hw-pipeline}

\textcolor{red}{Todas as arquiteturas otimizadas usam \textit{pipeline}. Sem \textit{pipeline}, as iterações de um laço rodam uma depois da outra \cite[p.~43]{kastner-hls}. Com \textit{pipeline}, uma iteração começa antes de a anterior terminar \cite[p.~45]{kastner-hls}. No começo, o \textit{pipeline} está vazio, e só os primeiros estágios trabalham. O tempo até todos os estágios ficarem ocupados é o tempo de enchimento do \textit{pipeline}. Depois dele, todos os estágios trabalham ao mesmo tempo, cada um com um dado diferente \cite[p.~114]{kastner-hls}.}
% ANTIGO (2026-11-06): depois do cite p. 43, "Na versão sem otimizações, o pipeline automático de
%   laços foi desligado (Tabela tab:hw-arquiteturas)." Foi para o primeiro parágrafo da subseção.
% Definição do enchimento: redação própria. O livro do Kastner não tem uma frase literal
%   sobre "pipeline fill" ou "startup". As passagens mais próximas: p. 114 ("Once the
%   pipeline is full, all four subfunctions are executing concurrently, but each one is
%   operating on different input data"), p. 45 (exemplo abaixo) e p. 146 (esvaziamento).

\textcolor{red}{No exemplo do livro, cada iteração de um laço leva 4 ciclos, e o laço tem 11 iterações, com II de 1 ciclo. A primeira iteração termina no ciclo 4, e a última, no ciclo 14 \cite[p.~45]{kastner-hls}. Ou seja, o laço leva $4 + (11-1) \cdot 1 = 14$ ciclos. No fim do laço, o \textit{pipeline} precisa esvaziar antes de o código seguinte começar. Por isso, a profundidade do \textit{pipeline} é um custo fixo, que pesa mais quando o laço tem poucas iterações \cite[p.~146]{kastner-hls}.}

\textcolor{red}{O mesmo vale para tarefas inteiras, como a inversão de uma matriz. Seja $L$ a latência de uma tarefa e $I$ o intervalo em regime (Seção~\ref{subsubsec:latencia-intervalo}). A primeira tarefa termina depois de $L$, e cada tarefa seguinte termina $I$ depois da anterior. Assim, $n$ tarefas seguidas terminam em}
\begin{equation}\label{eq:fim-tarefas}
\textcolor{red}{t_n = L + (n-1)\, I.}
\end{equation}
\textcolor{red}{O limite de ciclos da Tabela~\ref{tab:limite-ciclos} vale para o intervalo $I$, e não para a latência $L$. As subseções a seguir mostram como cada otimização reduz o intervalo.}

% ---------------------------------------------------------------------
\subsubsection{\texorpdfstring{\textcolor{red}{Pipeline de laços}}{Pipeline de laços}}\label{subsubsec:hw-pipeline-lacos}

\textcolor{red}{A arquitetura 9090 coloca a diretiva \texttt{\#pragma HLS pipeline II=1} nos laços internos das etapas do módulo (Figura~\ref{fig:modulo-inversao}). Só a etapa \texttt{D\_extraction} fica sem a diretiva. Essa diretiva pede ao HLS um II de 1 ciclo. Se as dependências ou os recursos não permitirem, o HLS usa um II maior \cite[p.~45]{kastner-hls}.}
% FONTE: cholesky_hardware_diretivas/{lu_decomposition, foward_substitution, inversion,
%   g_step, s_step}.cpp (pipeline II=1 nos laços internos); dv_step.cpp (D_extraction) sem pragma.

\textcolor{red}{A Figura~\ref{fig:hw-pipeline-lacos} mostra um laço da decomposição LU. Sem \textit{pipeline}, cada iteração espera a anterior terminar. Com II = 1, uma iteração nova começa a cada ciclo. Os quatro estágios só ficam ocupados no ciclo 4, quando a primeira iteração termina. Assim, o enchimento dura cerca de uma latência de iteração. O esvaziamento, no fim do laço, também. As etapas e os laços de uma matriz ainda rodam um depois do outro. Assim, a próxima inversão só começa depois que a anterior termina \cite[p.~114]{kastner-hls}. Nessa arquitetura, o intervalo entre inversões é igual à latência de uma inversão.}

\begin{figure}[!ht]
\caption{\textcolor{red}{Loop pipelining in the LU decomposition. The code is simplified.}}
\label{fig:hw-pipeline-lacos}
\vspace{4pt} % a classe zera o espaço entre a caption e a figura
\centering
\begin{minipage}[c]{0.40\textwidth}
\color{red}\fontsize{7}{8.6}\selectfont
\begin{verbatim}
for (line = column + 1; line < n; line++)
  for (k = column; k < n; k++) {
#pragma HLS pipeline II=1
    t = L[line][column] * U[column][k];
    U[line][k] = U[line][k] - t;
  }
\end{verbatim}
\end{minipage}\hfill
\begin{minipage}[c]{0.57\textwidth}
\centering
\begin{tikzpicture}[font=\fontsize{6.5}{7.8}\selectfont,
  every node/.append style={text=red},
  est/.style={draw=black!60, fill=#1, minimum width=0.50cm, minimum height=0.32cm, inner sep=0pt, line width=0.3pt},
  est/.default=white]
\def\u{0.50}\def\h{0.38}
\definecolor{e1}{HTML}{CDE2FB}\definecolor{e2}{HTML}{9EC5F4}\definecolor{e3}{HTML}{6FA6E8}\definecolor{e4}{HTML}{D9D9D9}
% estágios de uma iteração: Rd, Mul, Sub, Wr (latência de 4 ciclos)
\newcommand{\iter}[3]{% #1 = ciclo de início (0-index), #2 = linha, #3 = rótulo
  \node[est=e1] at (#1*\u+0.5*\u, -#2*\h) {\fontsize{5}{6}\selectfont Rd};
  \node[est=e2] at (#1*\u+1.5*\u, -#2*\h) {\fontsize{5}{6}\selectfont Mul};
  \node[est=e3] at (#1*\u+2.5*\u, -#2*\h) {\fontsize{5}{6}\selectfont Sub};
  \node[est=e4] at (#1*\u+3.5*\u, -#2*\h) {\fontsize{5}{6}\selectfont Wr};
  \node[left] at (-0.05, -#2*\h) {#3};}
% (a) sem pipeline
\node[anchor=west] at (-1.0, 0.62) {(a) Without pipeline};
\iter{0}{0}{it.~1}
\begin{scope}[shift={(4*\u,0)}]
  \node[est=e1] at (0.5*\u,0) {\fontsize{5}{6}\selectfont Rd};
  \node[est=e2] at (1.5*\u,0) {\fontsize{5}{6}\selectfont Mul};
  \node[est=e3] at (2.5*\u,0) {\fontsize{5}{6}\selectfont Sub};
  \node[est=e4] at (3.5*\u,0) {\fontsize{5}{6}\selectfont Wr};
\end{scope}
\node at (6*\u, 0.30) {it.~2};
\node[right] at (8*\u, 0) {$\cdots$ next iteration after 4 cycles};
% (b) com pipeline, II = 1, 6 iterações
\begin{scope}[shift={(0,-1.30)}]
\node[anchor=west] at (-1.0, 0.62) {(b) Pipeline, II = 1};
\foreach \i in {0,...,5} { \pgfmathtruncatemacro{\r}{\i+1} \iter{\i}{\i}{it.~\r} }
% eixo dos ciclos
\foreach \c in {1,...,9} \node[text=red!70!black] at (\c*\u-0.5*\u, -6*\h+0.02) {\fontsize{5}{6}\selectfont\c};
\node[left] at (-0.05, -6*\h+0.02) {cycle};
% II
\draw[red, |<->|, line width=0.4pt] (0,0.24) -- (\u,0.24) node[midway, above=-1pt] {II};
% enchimento, regime e esvaziamento
\draw[decorate, decoration={brace, mirror, amplitude=3pt}, red, line width=0.4pt] (0.02,-6*\h-0.15) -- (3*\u-0.02,-6*\h-0.15) node[midway, below=2pt, align=center] {fill};
\draw[decorate, decoration={brace, mirror, amplitude=3pt}, red, line width=0.4pt] (3*\u+0.02,-6*\h-0.15) -- (6*\u-0.02,-6*\h-0.15) node[midway, below=2pt, align=center] {full};
\draw[decorate, decoration={brace, mirror, amplitude=3pt}, red, line width=0.4pt] (6*\u+0.02,-6*\h-0.15) -- (9*\u-0.02,-6*\h-0.15) node[midway, below=2pt, align=center] {flush};
\draw[decorate, decoration={brace, amplitude=3pt}, red, line width=0.4pt] (9*\u+0.12,0.15) -- (9*\u+0.12,-5*\h-0.15) node[midway, right=3pt, align=left] {loop latency\\$4 + (6-1)\cdot 1 = 9$};
\end{scope}
\end{tikzpicture}
\end{minipage}
\vspace{4pt}
\legend{\textcolor{red}{Source: the author. Rd: read, Mul: multiply, Sub: subtract, Wr: write. Iteration latency of 4 cycles, as in the example of the book \cite[p.~45]{kastner-hls}.}}
\end{figure}

% ---------------------------------------------------------------------
\subsubsection{\texorpdfstring{\textcolor{red}{Dataflow entre etapas}}{Dataflow entre etapas}}\label{subsubsec:hw-dataflow}

\textcolor{red}{A arquitetura 4545 mantém o \textit{pipeline} dos laços e junta as funções do módulo em quatro etapas. A primeira é a decomposição LU. A segunda calcula a matriz $\mathbf F$ (\texttt{D\_extraction}, \texttt{D\_square\_root} e \texttt{G\_decomposition}). A terceira é a substituição direta, e a quarta, a inversão. A diretiva \texttt{\#pragma HLS dataflow}, na função de topo, transforma cada etapa em um processo. Assim, uma matriz nova entra na primeira etapa assim que a matriz anterior passa para a segunda \cite[p.~114]{kastner-hls}.}
% FONTE: cholesky_hardware_diretivas_4545/cholesky_top.cpp (etapa_lu, etapa_g, etapa_fs,
%   etapa_inv; dataflow em cholesky_inverse, linha 68).

\textcolor{red}{O \textit{dataflow} forma um \textit{pipeline} de tarefas, controlado pela troca de sinais entre os processos. O \textit{pipeline} de laços, por sua vez, age ciclo a ciclo \cite[pp.~114--115]{kastner-hls}. Entre dois processos, os dados passam por uma FIFO ou por um \textit{buffer} ping-pong. O \textit{buffer} ping-pong tem dois blocos: o processo que produz os dados escreve em um, enquanto o processo que os consome lê o outro \cite[p.~115]{kastner-hls}.}

\textcolor{red}{Cada etapa leva a diretiva \texttt{\#pragma HLS inline off}. Essa diretiva impede o HLS de copiar a função para dentro da função que a chama, e mantém a hierarquia do código \cite[p.~52]{kastner-hls}. Assim, cada etapa continua sendo uma função separada e vira um processo.}

\textcolor{red}{Na Figura~\ref{fig:hw-dataflow}, a latência de uma matriz é a soma das quatro etapas. Uma etapa só recebe a matriz seguinte depois de terminar a anterior. Por isso, o intervalo é a duração da etapa mais longa. Sem \textit{dataflow}, a matriz 2 só começaria no fim da matriz 1.}
% SEM FONTE: "intervalo = etapa mais longa" é lido da figura (não achei frase no Kastner;
%   a p. 97 fala de "weakest link", mas de uma porta de memória).

\begin{figure}[!ht]
\caption{\textcolor{red}{Dataflow between the stages of the inversion. The code is simplified, and the durations are not to scale.}}
\label{fig:hw-dataflow}
\vspace{4pt}
\centering
\begin{minipage}[c]{0.33\textwidth}
\color{red}\fontsize{7}{8.6}\selectfont
\begin{verbatim}
void cholesky_inverse(A, inverse) {
#pragma HLS dataflow
  etapa_lu(A, L, U);
  etapa_g(U, L, G);
  etapa_fs(G, G_inv);
  etapa_inv(G_inv, inverse);
}

static void etapa_fs(G, G_inv) {
#pragma HLS inline off
  forward_substitution(G, G_inv);
}
\end{verbatim}
\end{minipage}\hfill
\begin{minipage}[c]{0.64\textwidth}
\centering
\begin{tikzpicture}[font=\fontsize{6.5}{7.8}\selectfont,
  every node/.append style={text=red},
  tar/.style={draw=black!60, fill=#1, minimum height=0.34cm, inner sep=0pt, line width=0.3pt}]
\def\u{0.40}\def\h{0.46}
\definecolor{m1}{HTML}{CDE2FB}\definecolor{m2}{HTML}{9EC5F4}\definecolor{m3}{HTML}{D9D9D9}
% durações esquemáticas: LU = 4, G = 1, FS = 2, INV = 2 (soma 9, maior 4)
\newcommand{\caixa}[5]{% #1 linha, #2 início, #3 duração, #4 cor, #5 matriz
  \node[tar=#4, minimum width=#3*\u cm, anchor=west] at (#2*\u, -#1*\h) {#5};}
\foreach \l/\n in {0/etapa\_lu, 1/etapa\_g, 2/etapa\_fs, 3/etapa\_inv}
  \node[left, font=\fontsize{6.2}{7.4}\selectfont\ttfamily] at (-0.05, -\l*\h) {\n};
% matriz 1
\caixa{0}{0}{4}{m1}{1}\caixa{1}{4}{1}{m1}{1}\caixa{2}{5}{2}{m1}{1}\caixa{3}{7}{2}{m1}{1}
% matriz 2
\caixa{0}{4}{4}{m2}{2}\caixa{1}{8}{1}{m2}{2}\caixa{2}{9}{2}{m2}{2}\caixa{3}{11}{2}{m2}{2}
% matriz 3
\caixa{0}{8}{4}{m3}{3}\caixa{1}{12}{1}{m3}{3}\caixa{2}{13}{2}{m3}{3}\caixa{3}{15}{2}{m3}{3}
% eixo
\draw[black!60, ->, line width=0.4pt] (0,-3*\h-0.3) -- (17.6*\u,-3*\h-0.3) node[right] {time};
% intervalo
\draw[red, |<->|, line width=0.4pt] (0,0.30) -- (4*\u,0.30) node[midway, above=-1pt] {interval $I$ = longest stage};
% latência da matriz 1
\draw[red, |<->|, line width=0.4pt] (0,-3*\h-0.55) -- (9*\u,-3*\h-0.55) node[midway, below=-1pt] {latency $L$ = sum of the stages};
% sem dataflow
\draw[red, dashed, line width=0.4pt] (9*\u,0.18) -- (9*\u,-3*\h-0.3);
\node[align=left, anchor=south west] at (9*\u+0.05,0.05) {without dataflow,\\matrix 2 starts here};
% fim da matriz 3
\draw[red, |<->|, line width=0.4pt] (0,-3*\h-1.0) -- (17*\u,-3*\h-1.0) node[midway, below=-1pt] {matrix 3 ends at $L + 2I$ (Equation~\ref{eq:fim-tarefas})};
\end{tikzpicture}
\end{minipage}
\vspace{4pt}
\legend{\textcolor{red}{Source: the author. Numbers in the boxes: index of the matrix.}}
\end{figure}

% ---------------------------------------------------------------------
\subsubsection{\texorpdfstring{\textcolor{red}{Lotes de matrizes em vias}}{Lotes de matrizes em vias}}\label{subsubsec:hw-lotes}

\textcolor{red}{Nas arquiteturas 606 e 455, o módulo recebe várias matrizes $\mathbf G_c$ de uma vez, cada uma de um PRB. Esse conjunto é um lote. O lote tem 12 matrizes na 606 e 16 na 455. As matrizes de um lote não dependem umas das outras. As etapas do lote formam um \textit{dataflow}, como na 4545, e um lote novo entra a cada intervalo de lote. O intervalo por matriz é o intervalo do lote dividido pelo número de matrizes do lote, e é ele que deve atender o limite da Tabela~\ref{tab:limite-ciclos}. O número de lotes de um bloco é $N_{\text{PRB}}$ dividido pelo tamanho do lote, arredondado para cima. O último lote vai incompleto. Na Equação~\eqref{eq:fim-tarefas}, a tarefa passa a ser o lote.}
% FONTE: cholesky_hardware_diretivas_606/cholesky_top.cpp (comentário do topo; LOTE e
%   VIAS em cholesky.h; dataflow na linha 315 com lote_ler, lote_lu, lote_g, lote_fs,
%   lote_inv, lote_gravar). LOTE/VIAS: 12/6 (606), 16/8 (455).

\textcolor{red}{As matrizes do lote ficam em bancos de memória separados, com a diretiva \texttt{\#pragma HLS array\_partition cyclic} no índice da matriz. A partição cíclica distribui os elementos entre os bancos de forma alternada. Por exemplo, com fator 2, o vetor 1, 2, \ldots, 9 vira os vetores 1, 3, 5, 7, 9 e 2, 4, 6, 8 \cite[p.~88]{kastner-hls}. Com fator igual ao número de vias, a matriz $m$ fica no banco $m \bmod \text{vias}$. Uma via é o conjunto de operadores que trata um banco. A 606 tem 6 vias, e a 455, 8.}

\textcolor{red}{A Figura~\ref{fig:hw-lotes} mostra a 606. Em cada ciclo, o laço com \textit{pipeline} faz a mesma conta em 6 matrizes, uma em cada via. O lote é tratado em grupos de 6 matrizes, que se alternam a cada ciclo. Como as matrizes não dependem umas das outras, uma iteração não espera o resultado da anterior. As divisões também alternam entre as matrizes do lote, uma por ciclo. Assim, o divisor, que tem \textit{pipeline}, não espera o resultado da divisão anterior.}

\begin{figure}[!ht]
\caption{\textcolor{red}{Batch of 12 matrices in 6 lanes, in the 606 architecture. The code is simplified.}}
\label{fig:hw-lotes}
\vspace{4pt}
\centering
\begin{minipage}[c]{0.43\textwidth}
\color{red}\fontsize{7}{8.6}\selectfont
\begin{verbatim}
#pragma HLS array_partition variable=U \
                     cyclic factor=VIAS dim=1
for (k = column; k < n; k++)
  for (g = 0; g < GRUPOS; g++) {
#pragma HLS pipeline II=1
    for (p = 0; p < VIAS; p++) {
      m = g * VIAS + p;
      t = L[m][line][column] * U[m][column][k];
      U[m][line][k] = U[m][line][k] - t;
    }
  }
\end{verbatim}
\end{minipage}\hfill
\begin{minipage}[c]{0.54\textwidth}
\centering
\begin{tikzpicture}[font=\fontsize{6.5}{7.8}\selectfont,
  every node/.append style={text=red},
  mat/.style={draw=black!60, fill=#1, minimum width=0.62cm, minimum height=0.34cm, inner sep=0pt, line width=0.3pt},
  via/.style={draw=black!60, fill=white, minimum width=0.78cm, minimum height=0.34cm, inner sep=0pt, line width=0.3pt}]
\def\u{0.86}
\definecolor{g0}{HTML}{CDE2FB}\definecolor{g1}{HTML}{D9D9D9}
% bancos: matriz m no banco m mod 6
\foreach \p in {0,...,5} {
  \pgfmathtruncatemacro{\q}{\p+6}
  \node at (\p*\u, 0.40) {bank \p};
  \node[mat=g0] at (\p*\u, 0) {$m_{\p}$};
  \node[mat=g1] at (\p*\u, -0.38) {$m_{\q}$};
  \node[via] (v\p) at (\p*\u, -1.05) {lane \p};
  \draw[black!60, ->, line width=0.4pt] (\p*\u, -0.56) -- (v\p.north);
}
\node[left] at (-0.38, 0) {group 0};
\node[left] at (-0.38, -0.38) {group 1};
% ciclos do laço de atualização
\begin{scope}[shift={(0,-1.75)}]
\foreach \c/\g/\k/\cor in {0/0/k/g0, 1/1/k/g1, 2/0/k{+}1/g0, 3/1/k{+}1/g1} {
  \node[mat=\cor, minimum width=1.22cm, minimum height=0.52cm, align=center, font=\fontsize{6}{7}\selectfont] at (\c*1.26+0.61, 0) {group \g\\[-1pt]column $\k$};
}
\node[left] at (-0.05, 0) {update};
\node[right] at (5.0, 0) {$\cdots$};
\foreach \c in {0,...,11} {
  \pgfmathtruncatemacro{\gg}{\c<6 ? 0 : 1}
  \ifnum\gg=0 \def\cc{g0}\else\def\cc{g1}\fi
  \node[mat=\cc, minimum width=0.42cm] at (\c*0.42+0.21, -0.55) {\fontsize{5.5}{6}\selectfont\c};
}
\node[left] at (-0.05, -0.55) {division};
\draw[black!60, ->, line width=0.4pt] (0,-0.92) -- (5.3,-0.92) node[right] {cycle};
\end{scope}
\end{tikzpicture}
\end{minipage}
\vspace{4pt}
\legend{\textcolor{red}{Source: the author. Update: one cycle per box, 6 matrices at a time. Division: one matrix per cycle, in a single divider; the numbers are the indices $m$ of the matrices.}}
\end{figure}

% ---------------------------------------------------------------------
\subsubsection{\texorpdfstring{\textcolor{red}{Arranjo sistólico}}{Arranjo sistólico}}\label{subsubsec:hw-sistolico}

\textcolor{red}{As arquiteturas 152 e 114 também usam lotes, com 32 matrizes. A 152 tem 2 vias, e a 114 tem 4. Cada via tem o seu próprio divisor. Além disso, o laço que percorre uma linha ou uma soma é desenrolado em células, uma para cada índice $k$, lado a lado. Cada célula só troca dados com a vizinha. Projetos assim são chamados de arranjos sistólicos, e podem surgir quando um laço interno é desenrolado \cite[p.~180]{kastner-hls}.}
% FONTE: cholesky_hardware_diretivas_152_sistolico/cholesky_top.cpp (comentário do topo);
%   LOTE/VIAS: 32/2 (152), 32/4 (114).

\textcolor{red}{Na substituição direta e na inversão, a soma parcial passa de uma célula para a seguinte, na mesma ordem de somas da versão sem otimizações. Assim, cada ciclo calcula uma soma inteira. Na decomposição LU, cada célula atualiza um elemento da linha, e cada ciclo atualiza uma linha inteira. Para as células lerem todos os índices $k$ no mesmo ciclo, cada elemento da linha fica em um banco próprio, com \texttt{\#pragma HLS array\_partition complete} \cite[p.~41]{kastner-hls}.}

\textcolor{red}{Na decomposição LU e na substituição direta, cada passo usa as linhas gravadas no passo anterior. As matrizes do lote se alternam a cada ciclo. Assim, entre a gravação de uma linha e a sua leitura, passam tantos ciclos quanto o tamanho do passo. Os passos curtos recebem ciclos vazios, até um número mínimo de ciclos. Com isso, uma linha só é lida depois de sair do \textit{pipeline}. A diretiva \texttt{\#pragma HLS dependence ... dependent=false} informa isso ao HLS. Essa diretiva passa ao HLS uma informação que ele não consegue deduzir do código, e cabe ao projetista garantir que ela é verdadeira \cite[pp.~149--150]{kastner-hls}. Se o \textit{pipeline} ficar mais profundo que o passo, o resultado fica errado. A simulação da netlist, que compara as inversas com as do código C (Seção~\ref{subsec:hw-potencia}), confere isso. A Figura~\ref{fig:hw-sistolico} mostra as células e os ciclos vazios.}
% FONTE: comentário do topo de cholesky_top.cpp da 152 (CHOL_PROF_LU, CHOL_PROF_FS;
%   lu_passo, fs_passo; "Se a profundidade real do pipeline passar de CHOL_PROF_*, o
%   resultado do RTL fica errado; a co-simulação confere isso").

\begin{figure}[!ht]
\caption{\textcolor{red}{Systolic array of the forward substitution and minimum step of the LU decomposition. The code is simplified.}}
\label{fig:hw-sistolico}
\vspace{4pt}
\centering
\begin{minipage}[c]{0.42\textwidth}
\color{red}\fontsize{7}{8.6}\selectfont
\begin{verbatim}
#pragma HLS array_partition variable=Ua \
                     complete dim=3
for (it = 0; it < LU_ITER; it++) {
#pragma HLS pipeline II=1
#pragma HLS dependence variable=Ua \
          type=inter dependent=false
  if (idx < reais)       // else: empty
    for (p = 0; p < VIAS; p++)
      for (k = 0; k < n; k++)  // cells
        Ua[m][line][k] = row[k] - l*piv[k];
}
\end{verbatim}
\end{minipage}\hfill
\begin{minipage}[c]{0.55\textwidth}
\centering
\begin{tikzpicture}[font=\fontsize{6.5}{7.8}\selectfont,
  every node/.append style={text=red},
  cel/.style={draw=black!60, fill=#1, minimum width=0.62cm, minimum height=0.42cm, inner sep=0pt, line width=0.3pt},
  cel/.default={white},
  seta/.style={black!60, ->, line width=0.4pt}]
\definecolor{c0}{HTML}{CDE2FB}\definecolor{c1}{HTML}{D9D9D9}
% (a) células em cadeia: s_k = s_{k-1} + a_k b_k
\node[anchor=west] at (-0.9, 1.20) {(a) Cells of one sum};
\foreach \k/\x/\l in {0/0/0, 1/1.05/1, 2/2.10/2} {
  \node[cel=c0] (c\k) at (\x, 0) {cell \l};
  \draw[seta] (\x, 0.55) node[above=-2pt] {$a_{\l}, b_{\l}$} -- (c\k.north);
}
\node (dots) at (3.05, 0) {$\cdots$};
\node[cel=c0] (cn) at (4.0, 0) {cell $k$};
\draw[seta] (4.0, 0.55) node[above=-2pt] {$a_{k}, b_{k}$} -- (cn.north);
\draw[seta] (-0.65, 0) node[left] {$0$} -- (c0.west);
\draw[seta] (c0) -- (c1); \draw[seta] (c1) -- (c2); \draw[seta] (c2) -- (dots); \draw[seta] (dots) -- (cn);
\draw[seta] (cn.east) -- ++(0.45, 0) node[right] {sum};
% (b) passos da LU com ciclos vazios
\begin{scope}[shift={(0,-1.85)}]
\node[anchor=west] at (-0.9, 1.00) {(b) Steps of the LU decomposition};
\def\w{0.40}
\foreach \i/\g in {0/0, 1/1, 2/0, 3/1, 4/0, 5/1} {
  \ifnum\g=0 \def\cc{c0}\else\def\cc{c1}\fi
  \node[cel=\cc, minimum width=0.40cm, minimum height=0.34cm] at (\i*\w+0.5*\w, 0) {\fontsize{5}{6}\selectfont g\g};
}
\foreach \i/\g in {6/0, 7/1} {
  \ifnum\g=0 \def\cc{c0}\else\def\cc{c1}\fi
  \node[cel=\cc, minimum width=0.40cm, minimum height=0.34cm] at (\i*\w+0.5*\w, 0) {\fontsize{5}{6}\selectfont g\g};
}
\foreach \i in {8,...,11}
  \node[draw=black!50, fill=white, densely dotted, minimum width=0.40cm, minimum height=0.34cm, inner sep=0pt, line width=0.4pt] at (\i*\w+0.5*\w, 0) {\fontsize{5}{6}\selectfont --};
\draw[decorate, decoration={brace, amplitude=3pt}, red, line width=0.4pt] (0.02, 0.28) -- (6*\w-0.02, 0.28) node[midway, above=2pt] {step $c$};
\draw[decorate, decoration={brace, amplitude=3pt}, red, line width=0.4pt] (6*\w+0.02, 0.28) -- (12*\w-0.02, 0.28) node[midway, above=2pt] {step $c+1$ (short)};
\draw[decorate, decoration={brace, mirror, amplitude=3pt}, red, line width=0.4pt] (8*\w+0.02, -0.28) -- (12*\w-0.02, -0.28) node[midway, below=2pt] {empty cycles};
\draw[black!60, ->, line width=0.4pt] (0,-0.75) -- (12*\w+0.3,-0.75) node[right] {cycle};
\end{scope}
\end{tikzpicture}
\end{minipage}
\vspace{4pt}
\legend{\textcolor{red}{Source: the author. g0 and g1: groups of matrices of the batch, one per cycle. Two groups are drawn for clarity: the 152 has 16 groups, and the 114 has 8. In (b), the short step is filled with empty cycles up to a minimum length, so that a row is read only after it leaves the pipeline.}}
\end{figure}

% ---------------------------------------------------------------------
\subsection{\texorpdfstring{\textcolor{red}{Test Matrices}}{Test Matrices}}\label{subsec:hw-matrizes}
% ---------------------------------------------------------------------
\textcolor{red}{O hardware foi testado com as matrizes $\mathbf{G}_c$ que o BCD-MMSE inverte no simulador. Um script roda o simulador com a mesma semente e salva cada $\mathbf{G}_c$ antes da inversão. O simulador não muda. Assim, as matrizes são exatamente as usadas nas curvas de SER, com 1 iteração do BCD-MMSE.}
% EN (versão anterior, 2026-11-05): The hardware was tested with the matrices $\mathbf{G}_c$ that the BCD-MMSE inverts in the simulator. A script runs the simulator with the same seed and saves each $\mathbf{G}_c$ before the inversion. The simulator does not change. So the matrices are exactly the ones used for the SER curves, with 1 BCD-MMSE iteration.
% FONTE: AMOSTRAGEM.md, seções 1 e 2 (exportar_matrizes.py troca a função de inversão
%   só na memória do processo; matrizes recalculadas por fora coincidem até 1e-12).

\textcolor{red}{Como as condições de teste só mudam $N$ e a portadora (Seção~\ref{subsec:setup_condicoes}), as matrizes foram exportadas uma vez para cada par da Tabela~\ref{tab:condicoes_teste}. São 5 pares, pois TDLA30-5 e TDLA30-10 têm o mesmo par.}
% ANTIGO (2026-11-05, repetia Test Conditions e a Tabela tab:condicoes_teste):
% \textcolor{red}{As condições de teste mudam só dois valores nessas matrizes: o número de amostras de ruído $N$ e a frequência de portadora. Condições com o mesmo $N$ e a mesma portadora dão as mesmas matrizes. Por isso, as matrizes foram exportadas uma vez para cada par: 192 amostras a 3,5~GHz (TDLA30-5 e TDLA30-10), 192 a 28~GHz (TDLA30-75), 168 a 28~GHz (TDLA30-300), 132 a 3,5~GHz (TDLC300-100) e 24 a 3,5~GHz (TDLB100-400).}
% EN (versão anterior, 2026-11-05): The test conditions change only two values in these matrices: the number of noise samples $N$ and the carrier frequency. Test conditions with the same $N$ and the same carrier give the same matrices. So the matrices were exported once for each pair: 192 samples at 3.5~GHz (TDLA30-5 and TDLA30-10), 192 at 28~GHz (TDLA30-75), 168 at 28~GHz (TDLA30-300), 132 at 3.5~GHz (TDLC300-100) and 24 at 3.5~GHz (TDLB100-400).
% FONTE: AMOSTRAGEM.md, tabela das pastas (2_amostra_matrizes/amostras<N>_fc<GHz>GHz).

\textcolor{red}{A modulação não muda $\mathbf{G}_c = E_s\mathbf{H}_c\mathbf{H}_c^H + \hat{\mathbf{R}}_{cc}$. Em cada grupo, as configurações QPSK (C, D, G e H) diferem das configurações 16-QAM (A, B, E e F) só na modulação. Por isso, as configurações QPSK não foram exportadas.}
% ANTIGO (2026-11-06): "As configurações 1C, 1D, 1G e 1H diferem de 1A, 1B, 1E e 1F só na modulação."
%   Citava só o grupo 1, mas as QPSK dos grupos 2 e 3 também não foram exportadas
%   (indice.csv: 16x16 só 1A, 1B, 1E, 1F, 3A, 3B, 3E, 3F; 32x32 só 2A, 2B, 2E, 2F).
% EN (versão anterior, 2026-11-05): The modulation does not change $\mathbf{G}_c = E_s\mathbf{H}_c\mathbf{H}_c^H + \hat{\mathbf{R}}_{cc}$. Configurations 1C, 1D, 1G and 1H differ from 1A, 1B, 1E and 1F only in the modulation. Therefore, the QPSK configurations were not exported.
% FONTE: AMOSTRAGEM.md, seção 4 ("As configurações 5 a 8 ficam fora por padrão").
% CUIDADO: a distribuição de G_c é a mesma para a mesma SNR. No ponto de operação, a
%   SNR do QPSK é menor, e G_c fica maior (tab:gc-ponto-operacao). As matrizes do QPSK
%   aparecem nos pontos de SNR baixa das configurações 16-QAM.

\textcolor{red}{Exportar todas as matrizes não é prático. Para cada par, a amostra tem duas partes. A primeira sorteia 3 matrizes para cada configuração e ponto de SNR, para cobrir os pontos de operação da simulação. A segunda pega as 5 matrizes com os valores mais extremos de quatro métricas: a maior parte real ou imaginária, o maior elemento da inversa, o menor autovalor e o número de condição. Esses são os piores casos de faixa e de precisão para o formato de ponto fixo.}
% EN (versão anterior, 2026-11-05): Exporting every matrix is not practical. For each pair, the sample has two parts. The first part draws 3 matrices for each configuration and SNR point, to cover the operating points of the simulation. The second part takes the 5 matrices with the most extreme values of four metrics: the largest real or imaginary part, the largest element of the inverse, the smallest eigenvalue and the condition number. These are the worst cases of range and precision for the fixed-point format.
% FONTE: AMOSTRAGEM.md, seção 4 (critérios "cobertura" e extremos; semente 0) e seção 5
%   (264 matrizes 16x16 e 132 matrizes 32x32 de cobertura por par, mais até 20 extremos).

% ---------------------------------------------------------------------
\subsection{\texorpdfstring{\textcolor{red}{Power and Energy}}{Power and Energy}}\label{subsec:hw-potencia}
% ---------------------------------------------------------------------
\textcolor{red}{A potência foi estimada com a atividade de chaveamento do projeto roteado. Para cada condição de teste, um testbench inverte 20 matrizes da amostra, uma depois da outra. As 20 matrizes são espalhadas pelas configurações e sorteadas com uma semente fixa: configurações 1A, 1B, 1E e 1F no circuito $16 \times 16$, e 2A, 2B, 2E e 2F no circuito $32 \times 32$. O circuito $16 \times 16$ também atende o grupo 3, mas a sua potência foi estimada só com matrizes do grupo 1. Nas arquiteturas com lote, o testbench inverte lotes inteiros, e são pelo menos 3 lotes, para medir o intervalo entre lotes. As 20 matrizes são repetidas em ordem até completar os lotes. Assim, a 606 inverte 36 matrizes, a 455 inverte 48, e a 152 e a 114 invertem 96.}
% ANTIGO (2026-11-06): o parágrafo terminava em "... 2A, 2B, 2E e 2F no circuito $32 \times 32$."
% FONTE (lotes): tb_potencia.cpp da 606 (linhas 57-71: lotes = max(3, ceil(20/LOTE)); matriz
%   (b*LOTE+m) % 20); coluna "matrizes" de resultados_servidor/**/potencia.csv: 36 (606),
%   48 (455), 96 (152 e 114), 20 nas outras.
% EN (versão anterior, 2026-11-05): The power was estimated with the switching activity of the routed design. For each test condition, a testbench inverts 20 matrices from the sample, one after the other. The 20 matrices are spread over the configurations of the circuit and drawn with a fixed seed: configurations 1A--1H for the $16 \times 16$ circuit, and 2A, 2B, 2E and 2F for the $32 \times 32$ circuit.
% FONTE: rodar_potencia.py (sinteses/*_vivado/potencia/): "20 matrizes ... configurações
%   1 a 8 [16x16] / 9 a 12 [32x32], espalhadas pelas configurações e sorteadas com
%   semente fixa"; tb_potencia.cpp. Nas versões com lote, o testbench inverte lotes
%   inteiros (pelo menos 3), então o número de matrizes pode passar de 20
%   (potencia_versoes.py). CONFERIDO (2026-11-06): ver o FONTE (lotes) acima.
% CORRIGIDO (2026-11-06): o texto dizia "configurações 1A--1H". rodar_potencia.py sorteia
%   das configurações 1 a 8 (numeração antiga), mas as QPSK (5 a 8 = 1C, 1D, 1G, 1H) não
%   foram exportadas: o indice.csv do 16x16 só tem 1-4 (1A, 1E, 1B, 1F) e 13-16 (3A, 3E, 3B,
%   3F). Logo, só 1A, 1B, 1E e 1F entram. O grupo 3 está na amostra, mas não na potência
%   (decisão da autora: não rodar de novo).

\textcolor{red}{O Vivado gera uma netlist funcional do projeto roteado. A netlist é simulada no simulador do Vivado com o testbench, e a atividade de cada sinal é salva em um arquivo SAIF. As inversas calculadas pela netlist são comparadas, byte a byte, com as inversas do código C. Depois, o arquivo SAIF é lido no projeto roteado, e o Vivado informa a potência.}
% EN (versão anterior, 2026-11-05): Vivado writes a functional netlist of the routed design. The netlist is simulated in the Vivado simulator with the testbench, and the activity of every signal is saved in a SAIF file. The inverses computed by the netlist are compared, byte by byte, with the inverses of the C code. Then the SAIF file is read in the routed design, and Vivado reports the power.
% FONTE: potencia.py, passos 1 a 5 (write_verilog -mode funcsim; xsim com -debug typical;
%   read_saif; report_power). TODO FONTE: guia de análise de potência do Vivado (UG907).

\textcolor{red}{A versão sem otimizações foi simulada nas 5 condições de teste. Cada arquitetura otimizada foi simulada só na sua condição de teste alvo. A TDLA30-10 usa as mesmas matrizes da TDLA30-5.}
% EN (versão anterior, 2026-11-05): The version without optimizations was simulated in the 5 test conditions. Each optimized architecture was simulated only in its target test condition. TDLA30-10 uses the same matrices as TDLA30-5.
% FONTE: memória tcc-potencia-cenario-por-versao; pacote_servidor_formatos/LEIAME.md.
%   Justificativa: na versão sem otimização, a potência quase não mudou entre as
%   condições (ex.: 0,047 x 0,046 W no TDLB100-400). Isso é resultado.

\textcolor{red}{A energia por inversão é a potência dinâmica multiplicada pelo tempo entre duas inversões. Esse tempo é o intervalo em regime, em ciclos de clock, medido na simulação da netlist, multiplicado pelo período de 10~ns. Nas arquiteturas com lote, o intervalo medido é o do lote. Ele é dividido pelo número de matrizes do lote, para dar o intervalo por matriz (Seção~\ref{subsubsec:hw-lotes}).}
% ANTIGO (2026-11-06): o parágrafo terminava em "... multiplicado pelo período de 10~ns."
% EN (versão anterior, 2026-11-05): The energy per inversion is the dynamic power multiplied by the time between two inversions. This time is the steady-state interval, in clock cycles, measured in the netlist simulation, multiplied by the 10~ns period.
% FONTE: memória tcc-pacote-servidor-2023 (convenção da energia). Intervalo tirado do
%   xsim_<cenário>.log do servidor, não do ciclos_medios do potencia.csv, que inclui o
%   enchimento do pipeline do primeiro lote.


% =====================================================================
\section{\texorpdfstring{\textcolor{red}{Organização dos resultados}}{Organização dos resultados}}\label{sec:organizacao-resultados}
% =====================================================================


% ---------------------------------------------------------------------
% 1.6  ORGANIZAÇÃO DOS RESULTADOS                 -> fim da Metodologia
% ---------------------------------------------------------------------
\textcolor{red}{Os resultados são agrupados pelas características físicas da estação base: o tamanho do circuito de inversão ($16 \times 16$ ou $32 \times 32$), o número $M$ de antenas e o número $C$ de \textit{clusters}. Dentro de cada grupo, variam a modulação, $K$ e a IoTh. Assim, há três grupos: as configurações 1A--1H ($M = 128$, $C = 8$, circuito $16 \times 16$), 2A--2H ($M = 256$, $C = 8$, circuito $32 \times 32$) e 3A--3H ($M = 256$, $C = 16$, circuito $16 \times 16$).}
% ORIGEM: Análise bits.md, "hard-constraints" (M, C, tamanho do circuito) e
%   "soft-constraints" (modulação, K, IoTh). Os nomes G1-G3 não vão para a tese.

